% Options for packages loaded elsewhere
% Options for packages loaded elsewhere
\PassOptionsToPackage{unicode,hidelinks}{hyperref}
\PassOptionsToPackage{hyphens}{url}
\PassOptionsToPackage{dvipsnames,svgnames,x11names}{xcolor}
%
\documentclass[
  letterpaper,
]{scrreprt}
\usepackage{xcolor}
\usepackage{amsmath,amssymb}
\setcounter{secnumdepth}{-\maxdimen} % remove section numbering
\usepackage{iftex}
\ifPDFTeX
  \usepackage[T1]{fontenc}
  \usepackage[utf8]{inputenc}
  \usepackage{textcomp} % provide euro and other symbols
\else % if luatex or xetex
  \usepackage{unicode-math} % this also loads fontspec
  \defaultfontfeatures{Scale=MatchLowercase}
  \defaultfontfeatures[\rmfamily]{Ligatures=TeX,Scale=1}
\fi
\usepackage{lmodern}
\ifPDFTeX\else
  % xetex/luatex font selection
\fi
% Use upquote if available, for straight quotes in verbatim environments
\IfFileExists{upquote.sty}{\usepackage{upquote}}{}
\IfFileExists{microtype.sty}{% use microtype if available
  \usepackage[]{microtype}
  \UseMicrotypeSet[protrusion]{basicmath} % disable protrusion for tt fonts
}{}
\makeatletter
\@ifundefined{KOMAClassName}{% if non-KOMA class
  \IfFileExists{parskip.sty}{%
    \usepackage{parskip}
  }{% else
    \setlength{\parindent}{0pt}
    \setlength{\parskip}{6pt plus 2pt minus 1pt}}
}{% if KOMA class
  \KOMAoptions{parskip=half}}
\makeatother
% Make \paragraph and \subparagraph free-standing
\makeatletter
\ifx\paragraph\undefined\else
  \let\oldparagraph\paragraph
  \renewcommand{\paragraph}{
    \@ifstar
      \xxxParagraphStar
      \xxxParagraphNoStar
  }
  \newcommand{\xxxParagraphStar}[1]{\oldparagraph*{#1}\mbox{}}
  \newcommand{\xxxParagraphNoStar}[1]{\oldparagraph{#1}\mbox{}}
\fi
\ifx\subparagraph\undefined\else
  \let\oldsubparagraph\subparagraph
  \renewcommand{\subparagraph}{
    \@ifstar
      \xxxSubParagraphStar
      \xxxSubParagraphNoStar
  }
  \newcommand{\xxxSubParagraphStar}[1]{\oldsubparagraph*{#1}\mbox{}}
  \newcommand{\xxxSubParagraphNoStar}[1]{\oldsubparagraph{#1}\mbox{}}
\fi
\makeatother


\usepackage{longtable,booktabs,array}
\usepackage{calc} % for calculating minipage widths
% Correct order of tables after \paragraph or \subparagraph
\usepackage{etoolbox}
\makeatletter
\patchcmd\longtable{\par}{\if@noskipsec\mbox{}\fi\par}{}{}
\makeatother
% Allow footnotes in longtable head/foot
\IfFileExists{footnotehyper.sty}{\usepackage{footnotehyper}}{\usepackage{footnote}}
\makesavenoteenv{longtable}
\usepackage{graphicx}
\makeatletter
\newsavebox\pandoc@box
\newcommand*\pandocbounded[1]{% scales image to fit in text height/width
  \sbox\pandoc@box{#1}%
  \Gscale@div\@tempa{\textheight}{\dimexpr\ht\pandoc@box+\dp\pandoc@box\relax}%
  \Gscale@div\@tempb{\linewidth}{\wd\pandoc@box}%
  \ifdim\@tempb\p@<\@tempa\p@\let\@tempa\@tempb\fi% select the smaller of both
  \ifdim\@tempa\p@<\p@\scalebox{\@tempa}{\usebox\pandoc@box}%
  \else\usebox{\pandoc@box}%
  \fi%
}
% Set default figure placement to htbp
\def\fps@figure{htbp}
\makeatother





\setlength{\emergencystretch}{3em} % prevent overfull lines

\providecommand{\tightlist}{%
  \setlength{\itemsep}{0pt}\setlength{\parskip}{0pt}}



 


\usepackage[a4paper]{geometry}
\usepackage{ccicons}
\usepackage{academicons}
\definecolor{orcidlogocol}{HTML}{A6CE39}
\usepackage{ragged2e}
\usepackage{tikz}
\usepackage{fontspec}
\setsansfont{AtkinsonHyperlegibleNext}[
    Path=fonts/next/,
    Scale=0.9,
    Extension = .otf,
    UprightFont=*-Regular,
    BoldFont=*-Bold,
    ItalicFont=*-RegularItalic,
    BoldItalicFont=*-BoldItalic,
    ]
\setmainfont{AtkinsonHyperlegibleNext}[
    Path=fonts/next/,
    Scale=0.9,
    Extension = .otf,
    UprightFont=*-Regular,
    BoldFont=*-Bold,
    ItalicFont=*-RegularItalic,
    BoldItalicFont=*-BoldItalic,
    ]
\setmonofont{AtkinsonHyperlegibleMono}[
    Path=fonts/mono/,
    Scale=0.9,
    Extension = .otf,
    UprightFont=*-Regular,
    BoldFont=*-Bold,
    ItalicFont=*-RegularItalic,
    BoldItalicFont=*-BoldItalic,
    ]
\makeatletter
\@ifpackageloaded{bookmark}{}{\usepackage{bookmark}}
\makeatother
\makeatletter
\@ifpackageloaded{caption}{}{\usepackage{caption}}
\AtBeginDocument{%
\ifdefined\contentsname
  \renewcommand*\contentsname{Table of contents}
\else
  \newcommand\contentsname{Table of contents}
\fi
\ifdefined\listfigurename
  \renewcommand*\listfigurename{List of Figures}
\else
  \newcommand\listfigurename{List of Figures}
\fi
\ifdefined\listtablename
  \renewcommand*\listtablename{List of Tables}
\else
  \newcommand\listtablename{List of Tables}
\fi
\ifdefined\figurename
  \renewcommand*\figurename{Figure}
\else
  \newcommand\figurename{Figure}
\fi
\ifdefined\tablename
  \renewcommand*\tablename{Table}
\else
  \newcommand\tablename{Table}
\fi
}
\@ifpackageloaded{float}{}{\usepackage{float}}
\floatstyle{ruled}
\@ifundefined{c@chapter}{\newfloat{codelisting}{h}{lop}}{\newfloat{codelisting}{h}{lop}[chapter]}
\floatname{codelisting}{Listing}
\newcommand*\listoflistings{\listof{codelisting}{List of Listings}}
\makeatother
\makeatletter
\makeatother
\makeatletter
\@ifpackageloaded{caption}{}{\usepackage{caption}}
\@ifpackageloaded{subcaption}{}{\usepackage{subcaption}}
\makeatother
\usepackage{bookmark}
\IfFileExists{xurl.sty}{\usepackage{xurl}}{} % add URL line breaks if available
\urlstyle{same}
\hypersetup{
  pdftitle={Bibliometrics Quick Notes},
  pdfauthor={Dr.~Stephan Gauch},
  colorlinks=true,
  linkcolor={blue},
  filecolor={Maroon},
  citecolor={Blue},
  urlcolor={Blue},
  pdfcreator={LaTeX via pandoc}}


\title{Bibliometrics Quick Notes}
\author{Dr.~Stephan Gauch}
\date{}
\begin{document}
\begin{titlepage}
\tikz [remember picture, overlay] %
\node [shift={(1cm,-1cm)}] at (current page.north west) %
[anchor=north west] %
{\includegraphics[scale=.15]{images/KB_logo_black-light.jpg}};

  \vfill
  \centering
  {\Huge \bfseries KB Quick Notes \par}
  \vfill
  {\Large 
   Dr. Stephan Gauch \href{https://orcid.org/0000-0002-4715-5400}{\textcolor{orcidlogocol}{\aiOrcid}}\ \par
   {\large Robert K. Merton Zentrum für Wissenschaftsforschung \par
   Humboldt-Universität zu Berlin }}
  \vfill
  \justify
  \ccby \\ This work is licensed under a \href{https://creativecommons.org/licenses/by/4.0/deed.en}{Creative Commons Attribution 4.0 License}
\end{titlepage}

\renewcommand*\contentsname{Table of contents}
{
\hypersetup{linkcolor=}
\setcounter{tocdepth}{2}
\tableofcontents
}

\bookmarksetup{startatroot}

\chapter{What is all this?}\label{what-is-all-this}

In the legal professions, there is a concept that in German carries the
wonderful name ``Loseblattsammlung'' (aka looseleaf binder, often as
part of a subscription then called looseleaf service). You pay for the
joy of receiving bundles of notes on some legal topic or issue, that you
then may neatly file away in a special folder in order to be ``up to
date'' on all the relevant facts and concerns surrounding said topic. In
a way, what you are currently looking at is also some sort of a
looseleaf binder, in that it aims to inform you on things. Yet, in
contrast to being ``all up to date all the time,'' it is rather
something that should get you ``up to speed'' in that it gives you a
starting point for a subset of topics and concepts in bibliometrics.
Also, there is no real service to subscribe to. Rather, these notes are
a starting point of starting points.

\section{But\ldots{} why?}\label{but-why}

``Why not write a proper textbook?'' you might say. And yeah, sure,
thanks, great suggestion! I probably could give a number of excuses. But
rather than those, I give you this. Also, I thought it would be best to
have something to just jump in without prior training in order to set a
basis for an inclusive discourse about the concepts and ideas discussed.
Having recently screened the state of bibliometrics in use, I see a lot
of contributions that would probably not qualify for publication in a
bibliometric core journal. I don't think that is a problem in itself. I
also don't subscribe to the notion that authors of such contributions
are per se ``incompetent'' or just have to ``see the light'' or should
be prevented from trying new things. Not necessarily (only) for moral
reasons - the latent elitism should be obvious - but for pragmatic ones.
There is a lot of talk about being ``responsible'' when working with
metrics. That's true for sure, right? Yet, it also is also quite funny
in my book, since one probably would have a bit of a hard time to find
vocal advocates for a) the ``irresponsible'' use of metrics, or, for
that matter, b) ``meaningless'' metrics or c) measuring what is
``irrelevant''. So these pleas don't really do much on their own.
Anyway! Rather than thinking of commandments, I opted for this:
Providing a simple starting point to allow interested newcomers to join
the party.

\section{Who (I think) might enjoy reading
this?}\label{who-i-think-might-enjoy-reading-this}

Well, first, foremost and obviously, anyone who wants to know what
bibliometrics is about.

Obvious candidates are:

\begin{itemize}
\tightlist
\item
  Students not having access to an introductory bibliometrics course and
  want to start learning.
\item
  Students that do have access to an introductory bibliometrics course,
  but struggle with getting a grip on the basic concepts and how they
  relate to each other.
\item
  Practitioners and administrators planning to use bibliometrics in
  their work and need a starting point (sic!) on what's what.
\item
  Policymakers who want to engage in a discourse with the bibliometric
  community and connect on the level of language use, what concepts
  bibliometricians believe to be relevant and why they believe that.
\item
  Anyone who wants to better understand a bibliometric study.
\end{itemize}

\section{Where to start with the Quick
Notes?}\label{where-to-start-with-the-quick-notes}

Is there a structure to the individual notes? Do they have to be read in
a particular order? Yes, and no. No, you may read this in any order you
see fit. Yes, starting with the basics might be a good idea, but I guess
you would have figured that out on your own and would not need a cue for
that. Also yes, I tend to bracket bibliometrics into the following
territories and did so in the Quick Notes as well.

\textbf{Evaluative bibliometrics}, the territory of the vertical, where
the tideous repetitiveness lives. What is ``countable'', based on an
argument of shared characteristics, e.g.~one citation being just like
any other citation, is stacked up to find out who are the ``best of the
best of the really very, very bestest''. For better or worse as numerous
reform initiatives suggest.

\textbf{Explorative bibliometrics} is the territory of the horizontal.
The realm of mapping and charting the old and the new or finding some
path from one peculiarity to another. How do topics evolve? Who
cooperates with whom? These sorts of things. The Explorative and the
Evaluative are quite neighborly, since you can try to translate many
horizontal phenomena into vertical signifiers of ``value'' or
``quality''. A bibliometric mapping exercise can be both speak for
``success in community building'' as well as ``representation of a
community''. Hence, borders here are sometimes rather loosely defined.
Yet, both can have rather different ideas of what is ``relevant'' or
``adequate''.\footnote{Using a ``citation window'' means to count the
  incoming citations to a contribution of a pre-defined period of time
  in order to prevent unfair comparison of old and new publications.
  After all, the old ones had more time to get cited. Makes sense in
  evaluation, but probably less so when the aim is to explore or map
  something. The difference between a) ``We need to compare fairly and
  therefore have to control for unfair advantage of old publications!''
  does not play well with b) ``Why cut off a part of the past of a
  field, when your interest is in the past of a field?''. There is, of
  course, a Quick Note on Citation Windows.}

Both of these territories share that they work ``with'' databases in
some way or another.

\textbf{Curative bibliometrics} is the territory responsible for keeping
things in order. Pruning, weeding things out, naming all the things
properly, etc. In practical terms: Making sure that mapping and counting
can be done ``adequately'' and in a proper orderly fashion. Obviously,
this is less working with the database but rather ``on'' the databases.

Finally, \textbf{reflexive bibliometrics}, a territory that I consider
myself the most with nowadays, is a home for those who do not
necessarily(!) work with databases or not necessarily(!) work on
databases - usually they can do both, too - but try to understand how
working with and working on databases shapes what bibliometricians do,
what others do with what bibliometricians do, and what others,
e.g.~researchers, do when observing what other others do, e.g.~funders,
with things that bibliometricians do and how all of this plays together.
So, inhabitants are rather interested in the idea of working on the
profession of bibliometrics and performativity of bibliometrics,
i.e.~how measures or maps rather ``produce'' what they seek to ``show''.

Anyway. I digress, again. Let's move on.

\section{How the Quick Notes are
structured}\label{how-the-quick-notes-are-structured}

The notes themselves definitely have an internal structure. All
documents start with a ``What is this about'' section, giving a rough
idea what concept the note is about. Next, you will find some sort of a
situating passage. Why is that concept important in bibliometrics? Where
applicable, you will find a passage on ``How does this work?'' giving
some verbalized idea of the procedures involved.

The Quick Notes do NOT contain code. Preposterous! Where is the
tutorial? Isn't bibliometrics also a craft? Of course it is! And there
might be more to come. Eventually. I guess. Anyway!

The last section contains limitations and critique towards or at least
related to the concept. These are not the result of a careful screening
exercise but rather me channeling the field. Last but not least:
References. Another starting point. Obviously, that selection is
illustrative and not exemplary.

Another aspect that structures the Quick Notes is their length. They are
almost insultingly short given the glorious purpose an ambitious reader
might throw at them. Needless to say: Each is a rabbit hole. A
bottomless pit of sorts. Each and every one of the topics makes for a
quite wonderful fetish. Consider these notes a shovel, at maximum a very
tiny and short ladder. The absolute minimum that these should achieve
though is getting you, dear reader, to be able to discuss these issues
with those that are more versed and have already dug a deeper hole for
themselves. The absolute maximum is getting you hooked on one or more of
these fetishes.

\section{That's an awful lot of starting points and not much of a
conclusion}\label{thats-an-awful-lot-of-starting-points-and-not-much-of-a-conclusion}

Indeed. This is intentional and pretty much by design. These notes are
supposed to leave you, dear reader, a tad unfulfilled. So it's not an
all-inclusive trip but rather a badly drawn map on a napkin to instill
the recklessness of, dare I say, adventure. Also, you might have
realised that the Quick Notes don't have a DOI. This is not a
coincidence! You should, for multiple reasons, not reference these
notes. First of all, they are, for reasons argued above and didactic
purposes, shallow, dramatically short, on the brink of being
oversimplified. The aim is to have you understand the basics. So, rather
than referencing these notes, I would rather have them shared.

Some final remarks: Why did you not make this a proper Wiki? Well, all
of this is licensed CC-BY. Go forth and wreak havoc. Why did you not
cover {[}insert your favourite fetish here{]}? Great question! Thank
you! Can I contribute to this or write a Quick Note myself? Color me
delighted! For the moment, I guess you might just contact me about this:
\href{mailto:stephan.gauch@hu-berlin.de}{\nolinkurl{stephan.gauch@hu-berlin.de}}.

Anyway.

Enjoy the ride.

Berlin, 18 Jul 2025

\section{Acknowledgments}\label{acknowledgments}

I want to thank a lot of people. First of all, the Working Group
``Competence development'' of the
\href{https://bibliometrie.info/en/}{Competence Network Bibliometrics
(KB)}, who participated in the selection of and discussion on these
topics. I also would extend my thanks to Sophia Dörner, Beatrice
Yefimov, and Najko Jahn, who supported the finalisation and polishing of
these Quick Notes. All the remaining shortcomings that survived these
scrutinous efforts are probably, no, eventually, intentional but in any
case my own.

The development of these Quick Notes was partially funded by the Federal
Ministry of Research, Technology and Space (BMFTR) in terms of the
VaMoKo project (FKZ 16WIK2101D).

The material is free to use and re-use, licensed under the
\href{https://creativecommons.org/licenses/by/4.0/deed.en}{CC BY 4.0
License}.

\part{Basics}

\chapter{Bradford's Law of
Scattering}\label{bradfords-law-of-scattering}

\section{What is Bradford's Law
about?}\label{what-is-bradfords-law-about}

Bradford's Law, formulated by Samuel C. Bradford in 1934, is a pattern
often found in bibliometric studies, namely that in bibliometrics pretty
much everything is distributed \emph{log-normal}, i.e., among other
characteristics, heavily right-skewed. The law describes the scatter or
dispersion of scientific literature, that, if scientific journals are
arranged in order of decreasing productivity of articles on a given
subject, they may be divided into a nucleus of journals (so-called core
journals) more particularly devoted to the subject and several groups or
zones containing the same number of articles as the nucleus when the
journals are increasingly less representative to the subject as a whole,
i.e.~they contain less works related to a specific matter. Bradford's
law then predicts the number of journals within the nucleus and zones.
The relationship of size of the nucleus and the zones is, at least
according to the original paper by Bradford: \(1:b:b^2\). This may not
be universally applicable today due to changes in the scientific
publication landscape. Yet, in principle, the concept still holds up in
the way Bradford's Law is famously summarized as:

\begin{quote}
A small core of journals will account for the majority of significant
scientific papers.
\end{quote}

\section{Why is Bradford's Law
important?}\label{why-is-bradfords-law-important}

Bradford's Law has been important for understanding and managing the
expanse of scientific literature, assisting librarians and information
scientists in identifying influential journals in a particular field,
thus enabling allocation of resources for library collections. In the
context of practical applications it may therefore support resource
management in libraries by focusing on core journals. It also aids
researchers in targeting the most focused journals for their work, both
for reading and publication purposes. In this sense, it is also
important for both evaluative and exploratory bibliometrics. Especially
when aiming to evaluate or explore a thematic field, be it a discipline
or a topic, using, at least in part, a journal-based strategy,
Bradford's Law will highlight both the strength and the weakness of a
journal-based approach. First and foremost, Bradford's Law provides an
excellent argument to focus on core journals to produce an effective
query for delineating a field or topic as it helps in identifying the
most productive journals quickly and at the same time provides an
argument, why there are limits to journal-based strategies; especially
in the 2nd zone, the number of journals is squared compared to the 1st
zone and this 1st zone already contains a factorised number of journals
compared to the original core journals. The strengths in relying on
Bradford's Law to design effective queries are that a journal-based
strategy will produce a vast number of results very quickly by
identifying the core journals, perhaps together with domain experts on
the subject in question. This benefit is usually in part qualified by
publisher practices aiming to bundle journals into subscription packages
providing marginal reductions when choosing said packages over
individual subscriptions. Finally, Bradford's Law may be used to give a
rough prediction of a field's size. Following the general rule of
relationship between the 1st and 2nd zones, very rough assumptions can
be made about how \emph{large} a scientific field might be.

\section{Limitations}\label{limitations}

\subsubsection{\texorpdfstring{Sticking to Bradford's Law will make your
query more
\emph{mainstream}}{Sticking to Bradford's Law will make your query more mainstream}}\label{sticking-to-bradfords-law-will-make-your-query-more-mainstream}

The main weaknesses lie in coverage. Over-reliance on Bradford's Law
will produce results that may over-represent the \emph{mainstream} of a
topic or field. In a similar vein, optimizing search strategies gearing
toward Bradford's law may not be possible in highly emerging topics or
fields, or with high degrees of interdisciplinarity, where the
publication strategies of authors may be geared toward more generalized
outlets due to the lack of specialized journals not being available in
early stages.

\subsubsection{Bradford's Law might be driven by multidisciplinary
applications}\label{bradfords-law-might-be-driven-by-multidisciplinary-applications}

Another limitation of Bradford's Law lies in the subjectivity of both
article classification and journal selection. Especially in the
long-tail of the distribution a single article in a journal will
contribute to increasing the skewness of the overall distribution. In
the case of bibliometrics, which is increasingly being used as a basis
for tracing actual developments of fields or topics, e.g.~in the context
of systematic reviews, a large number of papers are \emph{applications}
of bibliometrics rather than bibliometric research in the stricter
sense.

\subsubsection{Bradford's Law is (in some sense) field
dependent}\label{bradfords-law-is-in-some-sense-field-dependent}

Bradford's Law may not apply uniformly across different subjects or
disciplines. How it diverges from the ideal constants and coefficients
may be dependent on a number of characteristics, i.e.~a field's age, the
application orientation of a field, tendencies for interdisciplinarity,
rhythms of scientific outputs, funding etc. The digital age and open
access movements have further altered the publication landscape and
access patterns, which may also further affect the applicability and
usefulness of Bradford's Law.

\subsubsection{Bradford's Law abstracts from
quality}\label{bradfords-law-abstracts-from-quality}

Bradford's law makes absolutely no claim about quality. Scholars might
jump to the conclusion that articles in the fringes, i.e.~the 2nd zone,
may always be of \emph{low quality}. Even though anecdotal evidence
might reflect this, especially in fields of high
application-orientation, this can not qualify as a general rule. One
prominent example might be the Hirsch-Index, a prominent measure of the
productivity and impact of a researcher in their field. The article
introducing the Hirsch-Index was published in Proceedings of the
National Academy of Sciences (PNAS). Definitely not one of the core
journals, an influential paper nonetheless.

\section{Further Reading}\label{further-reading}

Bailón-Moreno, R., Jurado-Alameda, E., Ruiz-Baños, R., \& Courtial, J.
P. (2005). Bibliometric laws: Empirical flaws of fit.
\emph{Scientometrics, 63}(2), 209--229.
\url{https://doi.org/10.1007/s11192-005-0211-5}

Bradford, S. C. (1934). Sources of information on specific subjects.
\emph{Engineering, 26}(4), 85--86.

Glänzel, W., \& Thijs, B. (2012). Using `core documents' for detecting
and labelling new emerging topics. \emph{Scientometrics, 91}(2),
399--416. \url{https://doi.org/10.1007/s11192-011-0591-7}

Mutschke, P., \& Mayr, P. (2015). Science models for search: A study on
combining scholarly information retrieval and scientometrics.
\emph{Scientometrics, 102}(3), 2323--2345.
\url{https://doi.org/10.1007/s11192-014-1485-2}

Nicolaisen, J., \& Hjørland, B. (2007). Practical potentials of
Bradford's law: A critical examination of the received view.
\emph{Journal of Documentation, 63}(3), 359--377.
\url{https://doi.org/10.1108/00220410710743298}

Shenton, A. K., \& Hay-Gibson, N. V. (2011). Bradford's Law and its
relevance to researchers. \emph{Education for Information, 27}(4),
217--230. \url{https://doi.org/10.3233/EFI-2009-0882}

\chapter{Citation Windows}\label{citation-windows}

\section{What is a Citation Window?}\label{what-is-a-citation-window}

A citation window refers to the specific time period during which
citations to a contribution are counted and analyzed. This window can
range from a few years to several decades. In practice, however, the
periods for citation windows typically range between 2 and 5 years.

\section{Why is it Important?}\label{why-is-it-important}

Using citation windows is important because they directly influence the
calculation of impact of a contribution. They achieve this by addressing
a specific notion of \emph{fairness} that relates to the comparison of
recent and old research. The first argument for citation windows is that
old articles had more time to accumulate citations and therefore have an
unfair advantage over recent ones. The second argument is that in order
for a scientific article to be referenced at all, it usually needs a
certain amount of time to get recognised in the publication landscape.
The rather idealized idea is that an article is being read, has some
relevance in some research process, the research itself is being
performed, the article is being referenced in a manuscript, which again
is submitted to a journal and goes through rounds of peer review until
finally, after publication, and some additional time to be included in a
bibliographic database, the original contribution we are interested in,
has attracted a citation. Obviously, different fields have different
speeds at which research is cited. For instance, in rapidly evolving
fields like information technology, shortening the citation window may
be more helpful, while in disciplines with longer research cycles, such
as history, longer windows may be necessary.

\section{How Does it Work?}\label{how-does-it-work}

The length of a citation window strongly depends on the task at hand,
the requirements toward immediacy or the research question as well as
the practices of the field in question. Citations are then collected and
analyzed only within this predefined period. For instance, a 5-year
citation window starting from the year of publication means that only
citations received within those five years are considered in the
analysis. This information is not always part of publicly available
data.

\section{Limitations}\label{limitations-1}

\subsubsection{Does not account for slow
burns}\label{does-not-account-for-slow-burns}

One of the primary issues with citation windows is the risk of
misrepresenting the impact of research. Short windows may not capture
the long-term influence of a contribution, especially in fields where
citations accumulate slowly. Such \emph{sleeping beauties} might receive
little attention at first but a surge in reception after some time. If
this period is no longer within the citation window this reception is
being ignored.

\subsubsection{The longer the citation window - the older (and less
relevant) the
results}\label{the-longer-the-citation-window---the-older-and-less-relevant-the-results}

Longer windows may include citations that are less relevant to the
current state of the field, due to different rhythms of scientific
output. In some fields publication propensity favors shorter time frames
from research to publishing. This effect may even be catalyzed further
when considering differences in project or research process length. In
fields where rhythms are tight, longer citation windows might cover
things that are no longer relevant. Yet, context matters,
i.e.~distinctions have to be made between citation windows as a means to
discount relevance vs.~a means of preventing measures from being overly
influenced by matters of temporality.

\subsubsection{Citation windows might not fit the notion to be evaluated
against}\label{citation-windows-might-not-fit-the-notion-to-be-evaluated-against}

Even though citation windows are useful for matters of comparison it may
not be sensible to use when temporality is explicitly part of the
concept to evaluate against. One example can be the Hirsch-Index, which
puts a strong premium on \emph{breadth of impact}, which in turn makes
it a bad measure to use in the context of evaluation of junior
scientists. In the case of the Hirsch-Index the notion of \emph{breadth}
does not match up to the measurement logic of restricting temporality.
Similarly other evaluations that somehow address notions of
\emph{long-termedness} will not benefit from citation windows.

\subsubsection{Citation windows might be counterproductive for
exploratory
purposes}\label{citation-windows-might-be-counterproductive-for-exploratory-purposes}

Citation windows make intuitive sense in the context of evaluative
bibliometrics. In exploratory bibliometrics they might be
counter-productive. For instance, when questions of evolution or
dynamics of a field are relevant to the purpose of exploring. More often
than not, the exploratory bibliometrics involve understanding the
\emph{origin story} of a field or topic. Where did it come from? How did
it emerge? In this case limiting analysis by using citation windows may
eliminate just these interesting classics, e.g.~from co-citation
analyses.

\subsubsection{Citation windows only address one
issue}\label{citation-windows-only-address-one-issue}

Time isn't the only factor that influences citation counts. Document
type is another very plausible candidate with review papers usually
receiving a considerable premium on citations flowing in. Moreover, the
choice of a citation window can introduce bias, as it might favor
certain types of publications or disciplines over others. This might not
be a strong limitation or argument against the use of citation windows
but should just be a reminder that citation windows do not fix all
potential biases.

\subsubsection{Citation windows may be inherently
inaccurate}\label{citation-windows-may-be-inherently-inaccurate}

When calculating citation windows analysts have to be aware of what
information specifically the citation window is calculated on.
Sometimes, when calculation is performed using the publication year,
publications from early months in the year may receive a substantial
premium. Using the actual publication date might change the story quite
a bit.

\subsubsection{Be aware of difference in citation
windows}\label{be-aware-of-difference-in-citation-windows}

There is no universal standard when it comes to citation window length.
When interpreting or merging data the mere information that a citation
window has been applied may therefore be insufficient. There's also the
challenge of comparing studies using different citation windows, which
can lead to inconsistencies in interpretation of indicators.

\section{Further Reading}\label{further-reading-1}

Campanario, J. M. (2011). Empirical study of journal impact factors
obtained using the classical two-year citation window versus a five-year
citation window. \emph{Scientometrics, 87}(1), 189--204.
\url{https://doi.org/10.1007/s11192-010-0334-1}

Donner, P. (2018). Effect of publication month on citation impact.
\emph{Journal of Informetrics, 12}(1), 330--343.
\url{https://doi.org/10.1016/j.joi.2018.01.012}

Glänzel, W. (2004). Towards a model for diachronous and synchronous
citation analyses. \emph{Scientometrics, 60}(3), 511--522.
\url{https://doi.org/10.1023/B:SCIE.0000034391.06240.2a}

Glänzel, W., Schlemmer, B., \& Thijs, B. (2003). Better late than never?
On the chance to become highly cited only beyond the standard
bibliometric time horizon. \emph{Scientometrics, 58}(3), 571--586.
\url{https://doi.org/10.1023/B:SCIE.0000006881.30700.ea}

Wang, J. (2013). Citation time window choice for research impact
evaluation. \emph{Scientometrics, 94}(3), 851--872.
\url{https://doi.org/10.1007/s11192-012-0775-9}

\chapter{Coverage in Bibliometrics}\label{coverage-in-bibliometrics}

\section{What does Coverage mean in
Bibliometrics?}\label{what-does-coverage-mean-in-bibliometrics}

In bibliometrics, \emph{coverage} refers to the extent to which a
bibliographic database (or tool) includes relevant publications,
journals, conferences, and other scholarly outputs within a specific
field or across multiple disciplines. Coverage can be evaluated in terms
of the breadth (range of subjects, disciplines, or publication types)
and depth (historical range, level of detail in indexing) of the
included materials. Effective coverage is crucial for comprehensive
bibliometric analyses, as it directly impacts the accuracy and
reliability of the results derived from the data. Besides the notion of
coverage from a disciplinary perspective, it can also refer to the
occurrence of missing data on the level of individual records in a
database, which may also be referred to as \emph{completeness}.

\section{Why is it Important?}\label{why-is-it-important-1}

There is no complete database of scientific publications including all
the features that would qualify current bibliometric analyses. In that
sense all databases are defined by a focus, some might say bias, towards
specific criteria of inclusion. Both notions of coverage have
substantial implications for bibliometric analyses. For instance,
adequate coverage ensures a more complete and accurate representation of
the scholarly landscape, which in turn enables more valid comparisons
across different fields, institutions, or time periods. Despite its
importance for bibliometrics, balancing quality, i.e., more precisely,
and more problematically, citation counts, vis-a-vis quantity or
completeness might not be as straightforward as one might expect in
matters of coverage. Among the current considerations are matters of
post-colonial concern, e.g.~contributions by the \emph{periphery} of the
science system such as the global south, with some scholars rejecting
the notion of a primacy of the center vs.~the periphery. Other aspects
that problematize coverage in this way is identification of trends and
gaps, where a higher coverage assists in identifying emerging research
trends and potential gaps in the literature. The issue therefore is more
complex than \emph{the more, the merrier}.

Evaluating and ensuring coverage usually involves multiple aspects, such
as:

\begin{enumerate}
\def\labelenumi{\arabic{enumi}.}
\tightlist
\item
  Assessment of Bibliometric Sources: Examining the scope and extent of
  bibliometric databases and tools to determine their coverage.
\item
  Transparent Data Inclusion Criteria: Defining criteria for what types
  of publications and sources are included.
\item
  Regular Updates and Expansion: Continuously updating and expanding
  databases to include new publications, journals, and other relevant
  scholarly outputs.
\item
  Cross-Database Comparison: Comparing data across multiple bibliometric
  sources to identify coverage overlaps and gaps.
\end{enumerate}

\section{How Does it Work?}\label{how-does-it-work-1}

Coverage is a concept that is more a matter of consideration than
something that can be \emph{performed} as a method in the stricter
sense. More often than not, due to the commercial nature of
bibliographic data and available resources. With databases from the
Web-of-Science-Family the criteria usually will be a derivative of
citations over time relative to publication volume on the level of
journals taking into account the disciplinary focus of said journals,
a.k.a \emph{quality}. Scopus seemingly aims for more of a balance
between quality and quantity. Yet, in both cases, how, when and why a
journal is included is not overly transparent. Moreover, which journals
are being \emph{put to review} is unclear as well, rendering the overall
composition a bit of a mystery at times. Descriptions of working groups
or committees in part exist, yet these are also only moderately
informative toward the issue. The prevalence of focus is not just
present in established databases but also extends to current contenders.
Sometimes at the cost of the second notion of coverage. The platform
\emph{Dimensions} features a larger collection of articles leading to an
overall larger citation network but has issues with completeness.

\section{Limitations}\label{limitations-2}

\subsubsection{Adequate Coverage is
pricey}\label{adequate-coverage-is-pricey}

Achieving and maintaining extensive coverage requires significant
resources and effort. This is especially true when understanding
coverage as a dynamic concept rather than a static silo logic at a
definite point in time. Main drivers are changes in the topic, field and
publication landscape.

\subsubsection{Coverage is politics by other
means}\label{coverage-is-politics-by-other-means}

Coverage can be biased towards certain languages, regions, or
disciplines, leading to skewed perspectives both from an evaluative as
well as explorative perspective. Keeping up with the rapidly evolving
landscape of scholarly publications is a continuous challenge.

\subsubsection{Coverage is exclusion}\label{coverage-is-exclusion}

Coverage should not be confused with unreflected inclusion. Selecting
journals, repositories, hubs or articles always involves a judgment,
even if done \emph{100\% algorithmically}. In this sense understanding
coverage can (and maybe should) be understood more along the lines of
curation. Extensive coverage may result in an overwhelming amount of
garbage data, complicating analysis and interpretation.

\section{Further Reading}\label{further-reading-2}

Daniel, B. Klein, \& Chiang, E. (2004). The Social Science Citation
Index: A Black Box---With an Ideological Bias? \emph{Econ Journal Watch,
1}(1), 134--165.

Harzing, A.-W., \& Alakangas, S. (2016). Google Scholar, Scopus and the
Web of Science: A longitudinal and cross-disciplinary comparison.
\emph{Scientometrics, 106}(2), 787--804.
\url{https://doi.org/10.1007/s11192-015-1798-9}

Larsen, P. O., \& Von Ins, M. (2010). The rate of growth in scientific
publication and the decline in coverage provided by Science Citation
Index. \emph{Scientometrics, 84}(3), 575--603.
\url{https://doi.org/10.1007/s11192-010-0202-z}

Mongeon, P., \& Paul-Hus, A. (2016). The journal coverage of Web of
Science and Scopus: A comparative analysis. \emph{Scientometrics,
106}(1), 213--228. \url{https://doi.org/10.1007/s11192-015-1765-5}

Stahlschmidt, S., \& Stephen, D. (2020). Comparison of Web of Science,
Scopus and Dimensions databases. Berlin: Deutsches Zentrum für
Hochschul- und Wissenschaftsforschung.
\url{https://bibliometrie.info/downloads/DZHW-Comparison-DIM-SCP-WOS.PDF}

\chapter{Institutional Disambiguation in
Bibliometrics}\label{institutional-disambiguation-in-bibliometrics}

\section{What does disambiguation mean in
Bibliometrics?}\label{what-does-disambiguation-mean-in-bibliometrics}

Disambiguation of affiliation data refers to the process of accurately
identifying and distinguishing the institutional affiliations of authors
in scholarly publications. This involves resolving ambiguities and
variations in the way institutions are named or represented in
publication data. For instance, an institution could be referred to by
different names or acronyms, or multiple departments within a single
institution could be listed separately. The aim is to ensure that each
publication is correctly attributed to the right institution. This also
includes diachronic changes, i.e.~affiliations splitting up or merging,
like in the case of the Karlsruhe Institute of Technology or (at least
temporarily) the Berlin Institute of Health. Apart from institutional
disambiguation, similar principles are being applied to disambiguate
author names, which usually proves to be a significantly harder
challenge for various reasons.

\section{Why is it Important?}\label{why-is-it-important-2}

Disambiguation is the prerequisite for accurate attribution. This aspect
might be the simplest factor why disambiguation is important in
bibliometrics: to ensure correct attribution of research output to
institutions, crucial for institutional rankings, reputation, and
funding. The notion of accuracy of attribution is not limited to
evaluative bibliometrics such as productivity assessments and impact
evaluations. Rather, it also is relevant, maybe even more relevant, in
exploratory bibliometrics, e.g.~in the context of analyses of
collaboration networks. Yet, disambiguation also provides further
benefits beyond bibliometrics, e.g.~in the context of in-house library
and resource management of organizations. Finally, disambiguation
supports overall research visibility of an organization by accurately
showcasing an institution's research contributions as well as enhancing
overall addressability, aiding in visibility and recognition in
industry, policy and the academic communities. All in all,
disambiguation provides the basis for fairer and more meaningful
comparisons between institutions and is indispensable for providing
clean and high quality data for policy-making and strategic decisions in
research management.

\section{How Does it Work?}\label{how-does-it-work-2}

There are numerous disambiguation approaches. Some will focus on the use
of structure-detection using multiple pieces of information. Others will
focus on fuzzy matching incorporating spelling errors. Other approaches
again might focus on so-called \emph{master list approaches} that are
basically extensively curated thesauri of spelling variants. Other
approaches again aim to use complementary data, such as WikiData, to
clean institution strings. In the Competence Network Bibliometrics the
current approach is to use regular expressions, basically matching for
delicately definable patterns rather than \emph{plain} searches, to
define an enormous set of positive and negative rules on how institution
strings are consolidated. These rules are applied to clean, if possible,
to the level of institutes and then aggregate back to the level of the
institution as a whole including start and end dates to account for
changes in institutional setups. Most of the approaches mentioned have a
quite similar workflow, which includes the following steps.

\begin{enumerate}
\def\labelenumi{\arabic{enumi}.}
\tightlist
\item
  Data Collection: Gathering affiliation data from consolidated data
  sources or publication records.
\item
  Identification of Variations: Recognizing different variations and
  representations of the same institution (differences based on the
  approach chosen).
\item
  Standardization and Matching: Standardizing the names and details of
  institutions and matching different variants to a single, standardized
  form.
\item
  Continuous Updating: Regularly updating the disambiguation process to
  accommodate new institutions, mergers, name changes, etc.
\end{enumerate}

The recursiveness and circular nature of this idealized approach should
make it clear that disambiguation usually is not a \emph{one-shot} task
but rather a continuous effort. Yet, for specific approaches and
questions using smaller datasets it may be admissible to use combined
approaches including structure detection and manual inference to arrive
at high quality data.

\section{Limitations}\label{limitations-3}

\subsubsection{All data is dirty all the
time}\label{all-data-is-dirty-all-the-time}

Disambiguation can quickly become a highly complex and resource
intensive task. Assuming that commercial, or non-commercial for that
matter, databases are sufficient for evaluatory or exploratory purposes
might be a gross lapse of judgment. Even though newer approaches aim to
mitigate the problem and some providers try to shift cleaning either
into algorithms, communities or clients, there currently is no database
that is perfectly clean in the regard described above, even though they
might claim so in sales pitches and promotional material.

\subsubsection{Disambiguating data is like feeding a
dragon}\label{disambiguating-data-is-like-feeding-a-dragon}

The process can be complex and resource-intensive, requiring
sophisticated algorithms and expert intervention. All these things are
costly and will require continuous and/or distributed efforts to
achieve. As soon as manual inference comes into play knowledge about the
national or local organizational landscape both within and beyond
stereotypical \emph{research organizations} is an absolute must! At
least for now high quality data simply won't come cheap. Also,
affiliation data is continually changing, making ongoing maintenance a
challenge as well as an imperative.

\subsubsection{Don't expect disambiguation to be standardized
(yet)}\label{dont-expect-disambiguation-to-be-standardized-yet}

Disambiguation will sometimes feature an obscene variability across
sources. How affiliation data is processed by the different data
providers is not consistent or comparable. This lack of standardisation
can make integration of data and harmonization efforts a complex task.

\section{Further Reading}\label{further-reading-3}

Daraio, C., Lenzerini, M., Leporelli, C., Naggar, P., Bonaccorsi, A., \&
Bartolucci, A. (2016). The advantages of an ontology-based data
management approach: Openness, interoperability and data quality.
\emph{Scientometrics, 108}(1), 441--455.
\url{https://doi.org/10.1007/s11192-016-1913-6}

Donner, P., Rimmert, C., \& Van Eck, N. J. (2020). Comparing
institutional-level bibliometric research performance indicator values
based on different affiliation disambiguation systems.
\emph{Quantitative Science Studies, 1}(1), 150--170.
\url{https://doi.org/10.1162/qss_a_00013}

Müller, M.-C., Reitz, F., \& Roy, N. (2017). Data sets for author name
disambiguation: An empirical analysis and a new resource.
\emph{Scientometrics, 111}(3), 1467--1500.
\url{https://doi.org/10.1007/s11192-017-2363-5}

Rimmert, C., Schwechheimer, H., \& Winterhager, M. (2017).
Disambiguation of author addresses in bibliometric databases - technical
report. Bielefeld: Universität Bielefeld, Institute for
Interdisciplinary Studies of Science (I²SoS).

Tang, L., \& Walsh, J. P. (2010). Bibliometric fingerprints: Name
disambiguation based on approximate structure equivalence of cognitive
maps. \emph{Scientometrics, 84}(3), 763--784.
\url{https://doi.org/10.1007/s11192-010-0196-6}

\chapter{Precision and Recall}\label{precision-and-recall}

\section{What are Precision and
Recall}\label{what-are-precision-and-recall}

One of the main findings of the Cranfield Experiments was the
development of two measures that help to formulize how improvement in
several areas of information science can be measured, e.g.~how good a
search query is, how well a classification algorithm corresponds to
expert judgment or how to evaluate the effectiveness of search systems.
In a way it shares characteristics with Alpha and Beta errors, which are
prominent measures in inferential statistics about asserting or
rejecting a Null hypothesis and whether this is right or wrong. To
understand Precision and Recall it is important to understand two
separate issues. First, if documents in a result set are assigned as
relevant or not relevant and second if that assignment is true or false,
resulting in 4 (relevant vs.~not relevant x true vs.~false) conditions.

\begin{quote}
True Positive (TP): A document is deemed relevant and this assessment is
true
\end{quote}

\begin{quote}
False Positive (FP): A document is deemed relevant and this assessment
is false
\end{quote}

\begin{quote}
True Negative (TN): A document is deemed irrelevant and this assessment
is true
\end{quote}

\begin{quote}
False Negative (TN): A document is deemed irrelevant and this assessment
is false
\end{quote}

\textbf{Precision} refers to the proportion of retrieved documents that
are relevant. This is calculated in the following way: Take the number
of relevant documents found (True Positives) and divide those by the
total number of documents retrieved (the sum of True Positives and False
Positives), or, in mathematical terms: \(\dfrac{TP}{(TP+FP)}\).

This can be understood as a measure of \emph{correctness}, i.e.~a
perfect Precision means:

\begin{quote}
All documents found are relevant!
\end{quote}

\textbf{Recall}, on the other hand, refers to the proportion of relevant
documents found. This is calculated in the following way: Take the
number of relevant documents found (True Positives again) and divide
those by all relevant documents there are (the sum of True Positives and
False Negatives), or, again in mathematical terms:
\(\dfrac{TP}{(TP+FN)}\).

This can be understood as a measure of \emph{completeness}, i.e.~a
perfect Recall means:

\begin{quote}
All relevant documents have been found!
\end{quote}

Precision and Recall sometimes are reported as percentages, e.g.~90\%,
or in decimal notation, 0.9.

\section{Why is it Important?}\label{why-is-it-important-3}

In bibliometrics, Precision and Recall are crucial for assessing the
performance of information retrieval systems, such as search engines and
digital libraries, the quality of classification algorithms, procedures
of disambiguation of authors or affiliations, quality of search queries
and much more. It also helps to balance optimization strategies when
doing bibliometric research. For instance, when the potential of manual
post-processing is given, e.g.~in an exploratory analysis when the
corpus of publications will be screened for publications fitting the
intended field or topic delineation, a \emph{high recall strategy} will
overall yield better results because the initial corpus will be larger
and irrelevant things can be weeded out.

A similar argument can be made for classification algorithms that are
being post-processed manually. Over-indexing on Recall of course could
be problematic, too, leading to a search strategy that simply returns,
in the most extreme cases, all records in a database. High precision
strategies may be more appropriate when defining a seed, e.g.~when
post-processing a search strategy by incorporating structural aspects
such as citation networks. In this case, optimizing for precision might
be worthwhile as moving back and forth in the citation tree will
automatically increase recall. Overdoing high-precision strategies is
not optimal either. In the most extreme case a single document might be
returned which is really close to the topic but several other documents
are not being found, making the seed highly dependent on this single
paper leaving out substantial components, e.g.~competing paradigms.

\section{Limitations}\label{limitations-4}

\subsubsection{You usually can't have it
all}\label{you-usually-cant-have-it-all}

There usually is a substantial trade-off between Precision and Recall.
Often, improving Precision leads to a decrease in Recall, and vice
versa, making it challenging, if not impossible, to optimize both to
very high values.

\subsubsection{Relevance does all the
work}\label{relevance-does-all-the-work}

Sometimes, especially when doing explorative work, assigning relevance
may be a very challenging task, especially when a bibliometrician has
little or no domain knowledge about the field to be analyzed.
Furthermore, relevance is at times elusive and highly context dependent
subjecting relevance, and hence, Precision and Recall, depending on the
user's need and the specific application.

\subsubsection{The recursive problem of
completeness}\label{the-recursive-problem-of-completeness}

When working exploratively a complete set of relevant documents IS the
actual goal. Yet, Recall implies that we know this in advance, framing
the goal to attain a complete set of relevant documents being in part
dependent on having information about the complete set of relevant
documents. This is obviously recursive and therefore problematic. It
also shows the limits of Precision and Recall, i.e.~carefully
constructed (sic!) sets of documents and assessment of relevance. In
many real-world scenarios, it's therefore difficult to know the total
number of relevant documents, complicating the calculation of Recall. It
is usually more helpful to understand Precision and Recall as
\emph{flavors of optimization} and general \emph{guideposts} on what to
achieve overall, namely correctness and completeness.

\section{Further Reading}\label{further-reading-4}

Baeza-Yates, R., \& Ribeiro-Neto, B. (2011). \emph{Modern information
retrieval: The concepts and technology behind search} (2nd ed.).
Addison-Wesley Publishing Company.

Cleverdon, C. W. (1972). On the inverse relationship of Recall and
Precision. \emph{Journal of Documentation, 28}(3), 195--201.
\url{https://doi.org/10.1108/eb026538}

Manning, C. D., Raghavan, P., \& Schütze, H. (2009). \emph{Introduction
to information retrieval}. Cambridge University Press.
\url{https://nlp.stanford.edu/IR-book/pdf/irbookonlinereading.pdf}

Powers, D. M. W. (2008). Evaluation: From precision, recall and
F-Measure to ROC, informedness, markedness and correlation. \emph{Mach.
Learn. Technol., 2}. \url{https://doi.org/10.48550/ARXIV.2010.16061}

\part{Evaluate}

\chapter{Composite Indicators in
Bibliometrics}\label{composite-indicators-in-bibliometrics}

\section{What is it?}\label{what-is-it}

In the field of bibliometrics, composite indicators (CIs) are
sophisticated metrics formulated by integrating multiple individual
indicators into a single, overarching index. In this sense CIs resemble
other statistical concepts found in general statistical applications,
e.g.~additive indices. Usually composite indicators are designed to
encapsulate a broader and more nuanced understanding of a complex
phenomenon, e.g research performance, impact, or productivity aiming to
integrate multiple facets of a phenomenon. They typically combine
diverse bibliometric dimensions such as citation counts, publication
volumes, journal impact factors, and collaboration metrics. In other
cases bibliometric metrics can be part of composite indicators. The
overall aim of CIs is to synthesize these varied elements to offer a
holistic measure that captures the complex nature of scholarly
activities.

\section{Why is it important?}\label{why-is-it-important-4}

The significance of composite indicators lies in their ability to
provide a holistic evaluation of a phenomenon. They overcome the
limitations of singular bibliometric measures by offering a balanced
perspective that encompasses various aspects of scholarly work. This
comprehensive approach is particularly attractive for policy-makers,
institutional administrators, research managers etc. as aid in informed
decision-making, strategic planning, and resource allocation as CIs
linearize a multi-dimensional concept into a singular metric. It is
something that both has the flair of \emph{holisticity} and
\emph{multidimensionality} and at the same time features the convenience
of being able to sort it (most of the time descending). Furthermore,
these indicators are instrumental in facilitating comparative analysis
across researchers, institutions, or countries, allowing for a more
equitable assessment by normalizing diverse bibliometric measures into a
unified framework. One example of a composite indicator in bibliometrics
is the Hirsch-Index, even though its computation is rather unique and it
does subsume composite indicators more in thought than as actual
reference to a type of metric. Another example for a composite indicator
are the different altmetric indicators combining different outlets of
(social) media into a single metric.

\section{How does it work?}\label{how-does-it-work-3}

The development and application of composite indicators in bibliometrics
involve a multi-step process that is intricately documented in the
\emph{OECD Handbook on Constructing Composite Indicators}. In a very
simplified way the construction follows steps.

Initially, relevant bibliometric indicators that reflect different
facets of research performance, or any other phenomenon for that matter,
are selected. These indicators are then polarity-harmonized (same
direction), normalized to ensure comparability, typically by converting
them to a common scale using a measure (e.g.~Z-Transformation, rescaling
etc\ldots) that limits the range to a predefined scale of possible
values or at least aims to produce a more uniform distribution. The next
steps involve the careful weighting and aggregation of these indicators,
with each component assigned a specific weight based on its perceived
importance in the overall assessment. Producing such weights can be very
challenging and depends in part on the outcome of the previous steps. In
principle it is an assessment of contribution to the overall phenomenon
to be measured, or, in more simple words, answering questions such as
``How many Tweets are worth a feature in a national newspaper?''.
Sometimes weighting is achieved by transforming all the contributing
metrics into a single \emph{currency}, e.g.~monetary values, even though
such a process is, at the very best, tedious and complex. The weighted
aggregation then results in the composite index, which then should be
subjected to rigorous validation to ensure its reliability and validity
both quantitatively, e.g.~by integrating it into inferential models, but
also semiotically by aiming to place it within the canon of previously
established indicators.

\section{Limitations}\label{limitations-5}

\subsubsection{Designing composite indicators involve careful evidence
and
processing}\label{designing-composite-indicators-involve-careful-evidence-and-processing}

The construction of these indicators is a complex task that requires
meticulous consideration in selecting components and determining their
respective weights. This process can introduce subjectivity,
particularly in the weighting stage, potentially skewing the final
outcome.

\subsubsection{Compressing is not
unfolding}\label{compressing-is-not-unfolding}

Over-reliance on a single composite score may also lead to
oversimplification, failing to capture the nuanced and multifaceted
nature of a phenomenon fully. Even though composite indicators are
argued to produce a holistic view of a subject, they in the end
linearize the mutli-dimensionality into a single vector of values. In
this regard, there are two fractions in bibliometrics. There are those
that embrace composite indicators for their merits and argue that at the
very least the overall values reflect a multi-dimensional phenomenon and
still provide the potential for rankings. Opposed to this perspective
are those that argue that such compression may help to produce simple
rankings but provide no guidelines for improvement as changes in the CIs
are not reflected by underlying changes of the corresponding dimensions.
At the very least these changes are overshadowed by the methodology. The
proponents of this position favor capturing multidimensionality through
visualization techniques, e.g.~by using radar or spider plots,
sacrificing the benefits of simple ranking.

\subsubsection{Garbage in - composited garbage
out}\label{garbage-in---composited-garbage-out}

Accuracy and reliability of composite indicators are heavily dependent
on the quality and availability of the underlying data.

\section{Further Reading}\label{further-reading-5}

El Gibari, S., Gómez, T., \& Ruiz, F. (2022). Combining reference point
based composite indicators with data envelopment analysis: Application
to the assessment of universities. \emph{Scientometrics, 127}(8),
4363--4395. \url{https://doi.org/10.1007/s11192-022-04436-0}

Johnes, J. (2018). University rankings: What do they really show?
\emph{Scientometrics, 115}(1), 585--606.
\url{https://doi.org/10.1007/s11192-018-2666-1}

Makkonen, T., \& Van Der Have, R. P. (2013). Benchmarking regional
innovative performance: Composite measures and direct innovation counts.
\emph{Scientometrics, 94}(1), 247--262.
\url{https://doi.org/10.1007/s11192-012-0753-2}

Moon, H. S., \& Lee, J. D. (2005). A fuzzy set theory approach to
national composite S\&T indices. \emph{Scientometrics, 64}(1), 67--83.
\url{https://doi.org/10.1007/s11192-005-0238-7}

Nasir, A., Ali, T. M., Shahdin, S., \& Rahman, T. U. (2011). Technology
achievement index 2009: Ranking and comparative study of nations.
\emph{Scientometrics, 87}(1), 41--62.
\url{https://doi.org/10.1007/s11192-010-0285-6}

OECD/European Union/EC-JRC. (2008). \emph{Handbook on constructing
composite indicators: Methodology and user guide}. Paris: OECD
Publishing. \url{https://doi.org/10.1787/9789264043466-en}

Vinkler, P. (2006). Composite scientometric indicators for evaluating
publications of research institutes. \emph{Scientometrics, 68}(3),
629--642. \url{https://doi.org/10.1007/s11192-006-0123-z}

\chapter{Field Normalization}\label{field-normalization}

\section{What is Field
Normalization?}\label{what-is-field-normalization}

Field normalization is a process used in bibliometrics to adjust for
variations across different scientific fields when comparing
bibliometric indicators, such as citation counts. This is necessary
because different scientific fields have distinct publication and
citation practices.

\section{Why is it important?}\label{why-is-it-important-5}

Imagine comparing a biology paper and a mathematics paper based solely
on their citation counts. This might be misleading because biology
papers, in general, tend to receive more citations due to the nature of
the field. Field normalization aims to level the playing field, allowing
for a fairer comparison.

\section{How does it work?}\label{how-does-it-work-4}

Determining expected values: Normalized indicators operate on a notion
of the \emph{average}, the \emph{usual}, the \emph{expected} etc\ldots{}
In simple terms: What we can use to assess if something is \emph{a lot}.
In order to normalize indicators, we need to calculate a baseline for a
specific field. This usually involves calculating the average number of
citations for papers in that field, published in a particular year. The
median, despite being more robust toward the outliers and right-skewed
distributions we find in citation counts, is often not a good measure as
the median often does not vary. More often than not, the median might
simply be 0. Other percentiles such as the Top 10\% or Top 1\%, i,e, the
90th-percentile or 99th-percentile, are used in another context.

Accounting for the expected values: In the next step, the citation count
of an individual paper is compared to this baseline. We achieve this by
simply dividing the observed value by the expected value calculated in
step 1. For example, if the average citation count in field X is 10 and
a paper in field X has 20 citations, it is performing above the average
for its field. Most normalized indicators will therefore be centered
around 1 as the \emph{neutral} outcome of normalization. The indicator
designer might choose to center an indicator at 0. In this case usually
a) 1 is subtracted from the ratio or b) a logarithm is applied to the
ratio. In the latter case, there is also an element of diminishing
returns built into the indicator.

\section{Limitations}\label{limitations-6}

\subsubsection{Not sensitive enough}\label{not-sensitive-enough}

Some critics of this approach argue that normalized indicators as they
are described above are not useful as they do not vary over time.

\subsubsection{The distribution of citations is too
skewed}\label{the-distribution-of-citations-is-too-skewed}

The arithmetic mean is sensitive to outliers. In the case of the
distribution of citation counts this is usually skewed to the right,
i.e.~there are a lot of papers rarely being cited and a few that are
cited disproportionately often. Critics argue that this makes the mean a
bad measure in such cases. They propose to use different measures such
as the so-called \emph{Excellence rate} or ppTOPX (ppTop10,
ppTop5\ldots) as it a) focuses on the politically relevant issue of
\emph{excellence} and ignoring \emph{noise} and b) is not prone to
distortion by the mean.

\subsubsection{(Often) too reliant on
classifications}\label{often-too-reliant-on-classifications}

This critique is mostly oriented toward field- or category-oriented
normalization. It does not account for normalization based on average
citation counts of journals such as the JCR (Journal Citation Rate). The
main point is that classifications may mis-represent what is happening
in science and underestimate the impact of inter- and transdisciplinary
research or at the very least over-represent what is happening in
\emph{core disciplines}. Among other limitations the Source Normalized
Citation Impact (SNIP) by Moed \& Colleagues addresses the problem of
classification by constructing the \emph{field} as a set of journals
citing a journal in question, thereby contextualizing the journal by its
reception.

\subsubsection{Practices of referencing can be biased against new
entrants}\label{practices-of-referencing-can-be-biased-against-new-entrants}

Some more fundamental criticism is put forward by critics relating to
the Matthew Effect. In a nutshell, the Matthew effect describes the
phenomenon that prestigious authors receive a premium on their
reputation, i.e.~they get more citations simply because they are more
well known and have a higher status or established reputation. This
argument usually calls into question the validity of evaluatory citation
analysis in general and is an issue of theories of citation.

\subsubsection{Citations reflect simplified notions of
impact.}\label{citations-reflect-simplified-notions-of-impact.}

A similar critique to the aforementioned argument calls into question
the validity of citations as measures of quality. This critique is
oriented toward the notion that impact should not be reduced to what is
often referred to as \emph{simplified counting}. At best, multiple
indicators should be used.

\subsubsection{Bibliometrics and Scientometrics have no theory of
scientific
fields}\label{bibliometrics-and-scientometrics-have-no-theory-of-scientific-fields}

Some critics argue that the scientometrics community has no convincing
theory of what delineates topics from scientific fields. Hence,
normalization efforts reflect mere categorization efforts of journals or
scientific articles that reflect rather artificial boundaries and
neither account for emerging scientific fields nor for phenomena of
interrelatedness such as inter- and transdisciplinarity.

\section{Further Reading}\label{further-reading-6}

Bornmann, L., Haunschild, R., \& Mutz, R. (2020). Should citations be
field-normalized in evaluative bibliometrics? An empirical analysis
based on propensity score matching. \emph{Journal of Informetrics,
14}(4), 101098. \url{https://doi.org/10.1016/j.joi.2020.101098}

Bornmann, L. (2020). How can citation impact in bibliometrics be
normalized? A new approach combining citing-side normalization and
citation percentiles. \emph{Quantitative Science Studies, 1}(4),
1553--1569. \url{https://doi.org/10.1162/qss_a_00089}

Heintz, B. (2010). Numerische Differenz. Überlegungen zu einer
Soziologie des (quantitativen) Vergleichs / Numerical difference. Toward
a sociology of (quantitative) comparisons. \emph{Zeitschrift für
Soziologie, 39}(3), 162--181.
\url{https://doi.org/10.1515/zfsoz-2010-0301}

Hicks, D., Wouters, P., Waltman, L., de Rijcke, S., \& Rafols, I.
(2015). Bibliometrics: The Leiden Manifesto for Research Metrics.
\emph{Nature, 520}(7548), 429--431.
\url{https://doi.org/10.1038/520429a}

Merton, R. K. (1968). The Matthew effect in science: The reward and
communication systems of science are considered. \emph{Science,
159}(3810), 56--63. \url{https://doi.org/10.1126/science.159.3810.56}

Moed, H. F. (2010). Measuring contextual citation impact of scientific
journals. \emph{Journal of Informetrics, 4}(3), 265--277.
\url{https://doi.org/10.1016/j.joi.2010.01.002}

\chapter{Fractional Counting}\label{fractional-counting}

\section{What is Fractional
Counting?}\label{what-is-fractional-counting}

Fractional counting relates to one of the major branches of evaluative
bibliometrics, i.e.~the evaluation of productivity or output volume. It
is a method used to attribute credit for scholarly publications,
particularly those with multiple authors. The overall goal is to account
for a notion of distributive fairness matching efforts.

\section{Why is it important?}\label{why-is-it-important-6}

In many fields, collaborative research leading to multi-author papers is
common. If the assessment of productivity is aimed to be relative to the
efforts involved, doubts about the validity of simply summing up the
number of contributions may emerge. Shouldn't it make a difference how
many authors are involved then? This is where fractional counting comes
in. The basic idea of fractional counting is to account for differences
in distribution of labor, i.e.~accounting for authors sharing the work
when compiling a scientific contribution. Assigning full credit to each
author can inflate the apparent productivity and impact of individuals
or institutions involved. Unlike full counting methods, where each
author or institution gets full credit for a publication, fractional
counting divides the credit among the contributors based on their
number.

\section{How Does It Work?}\label{how-does-it-work-5}

\subsubsection{Credit Allocation on different
levels}\label{credit-allocation-on-different-levels}

In a paper with multiple authors, instead of each author receiving one
full count for the publication, the credit is fractionally divided. For
example, if a paper has four authors, each author might receive 0.25
(1/4th) of a count for that paper. Yet, the level of allocation as well
as the level of analysis are important. Depending on the notion of what
\emph{contributes} to the paper the level of fractioning might vary.
Despite author-level fractioning, other levels may be used, e.g.~country
level, organization level, address level, i.e.~address of the
institution, which might be relevant for institutions with multiple
sites or locations, etc\ldots{} Specific forms may also involve pooling
the effort to one author, e.g.~the principal investigator, last author
position in medicine, or the first author representing the \emph{main
author}. Yet, these follow a completely different understanding of
fairness.

\subsubsection{Impact on Metrics}\label{impact-on-metrics}

Fractional counting affects various bibliometric indicators. This impact
is strictly not limited to analyses of publication volume and
productivity but also influences the way citation analyses are being
conducted.

\subsubsection{Variations in Practice}\label{variations-in-practice}

The exact method and data source of fractional counting can vary.
Comparing data from different sources, e.g.~WoS, Socpus, OpenAlex or
Google Scholar, as well as different levels of fractioning is strictly
not advised, let alone, merging them into one unified dataset for
analysis.

\section{Limitations}\label{limitations-7}

\subsubsection{Simple fractioning is not accounting for divisions of
labor}\label{simple-fractioning-is-not-accounting-for-divisions-of-labor}

Fractioning rests on the rather plausible idea that productivity is
relative to the number of authors involved. Yet, dividing by the number
of entities might not represent the actual efforts each author or
organization has invested in the paper. Recent developments in the
publication landscape try to account for this by introducing
contribution sections. In these sections, the authors state in what
parts of the research and publication process they have been involved
in.

\subsubsection{Fractional counting can put a penalty on
cooperation}\label{fractional-counting-can-put-a-penalty-on-cooperation}

Using fractional counting introduces a penalty on cooperation. This
might lead to misrepresentation of outputs. Imagine an organization that
a couple of years ago made it their mission to cooperate more with other
organizations. If that mission is being put to work, it is rather
plausible that the number of co-authored papers with authors from other
organizations will increase. This could lead to a situation where
full-counts, i.e.~whole counting of papers, increase but the number of
fractional papers level out as more authors are involved due to
cooperation activity. In extreme cases the fractional count could be
declining while full-count increases. Always use full counts as a litmus
test to understand your fractional counts.

\section{Further Reading}\label{further-reading-7}

Aksnes, D. W., Schneider, J. W., \& Gunnarsson, M. (2012). Ranking
national research systems by citation indicators. A comparative analysis
using whole and fractionalised counting methods. \emph{Journal of
Informetrics, 6}(1), 36--43.
\url{https://doi.org/10.1016/j.joi.2011.08.002}

Cronin, B. (2001). Hyperauthorship: A postmodern perversion or evidence
of a structural shift in scholarly communication practices?
\emph{Journal of the American Society for Information Science and
Technology, 52}(7), 558--69. \url{https://doi.org/10.1002/asi.1097}

Donner, P. (2024). Remarks on modified fractional counting.
\emph{Journal of Informetrics, 18}(4), 101585.
\url{https://doi.org/10.1016/j.joi.2024.101585}

Gauffriau, M., Larsen, P. O., Maye, I., Roulin-Perriard, A., \& Von Ins,
M. (2008). Comparisons of results of publication counting using
different methods. \emph{Scientometrics, 77}(1), 147--76.
\url{https://doi.org/10.1007/s11192-007-1934-2}

Leydesdorff, L., \& Opthof, T. (2010). Normalization at the field level:
Fractional counting of citations.
\url{https://doi.org/10.48550/ARXIV.1006.2896}

Sivertsen, G., Rousseau, R., \& Zhang, L. (2019). Measuring scientific
contributions with modified fractional counting. \emph{Journal of
Informetrics, 13}(2), 679--94.
\url{https://doi.org/10.1016/j.joi.2019.03.010}

Waltman, L., \& Van Eck, N. J. (2015). Field-normalized citation impact
indicators and the choice of an appropriate counting method.
\emph{Journal of Informetrics, 9}(4), 872--94.
\url{https://doi.org/10.1016/j.joi.2015.08.001}

\chapter{Hirsch-Index (h-index)}\label{hirsch-index-h-index}

\section{What is the Hirsch-Index?}\label{what-is-the-hirsch-index}

The Hirsch-Index, or h-index, is a metric developed by physicist Jorge
E. Hirsch in 2005 to quantify the combined impact and productivity of a
scholar's published work. In contrast to other evaluative metrics
focusing either on productivity or reception of scientific publications
by balancing the quantity (number of publications) and quality
(citations) of scientific output seeking to combine both perspectives
giving an overall notion of the \emph{eminence} of a scholar or research
organization. The h-index is defined as the number of papers (h) that
have been cited at least h times. For instance, if a researcher has an
h-index of 10, it means they have published 10 papers that have each
been cited at least 10 times. The set of publications for which these
conditions are true, i.e.~they definingly contribute to the h-Index due
to their received number of citations, is sometimes labeled as the
\emph{h-core} and is sometimes analyzed separately.

\section{Why is it Important?}\label{why-is-it-important-7}

The h-index is important because it provides a single figure that can
suggest the significance and broad impact of a researcher's cumulative
research contributions. In that respect it is a rather small scale
version of a composite indicator. It's widely used in academia for
evaluating the influence of a researcher's work, considering both the
quantity and quality of their publications. It's also employed by
funding agencies, universities, and research institutions for hiring,
promotion, and funding decisions. Its composition makes it resistant to
outliers, as e.g.~one or two highly cited papers, or a large number of
scarcely cited papers, have little effect on the Hirsch-Index. What
counts is continuous performance above a certain level. This also means
that with increase of the Hirsch-Index raising the index by 1 involves
increasingly higher efforts if trying to \emph{optimize} it. Finally, it
is exceedingly easy to compute, making it a measure that can be
performed with considerably little data as no field normalization is
involved. Another benefit of the Hirsch-Index is its scalability to
higher levels of aggregation making it a measure that is also applicable
on the level of organizations. Despite its prominence and widespread
application it is also an object of strong criticism in the current
evaluation landscape.

\section{How Does it Work?}\label{how-does-it-work-6}

Calculating the h-index involves the following steps:

\begin{enumerate}
\def\labelenumi{\arabic{enumi}.}
\tightlist
\item
  Listing Publications and Citations: Compile a list of the researcher's
  publications, along with the number of citations each paper has
  received.
\item
  Ranking by Citation Count: Rank the publications according to the
  number of citations.
\item
  Find the h-index threshold: Find the point where the highest rank of
  the paper in the list is equal to or greater than the number of
  citations it has received. This number is the h-index.
\end{enumerate}

\section{Limitations}\label{limitations-8}

\subsubsection{The Hirsch-Index favors senior
researchers}\label{the-hirsch-index-favors-senior-researchers}

The Hirsch-Index is only meaningfully interpretable when used to
understand differences within the group of established researchers for
the simple fact that the output is the first limiting factor for the
Hirsch-Index. It needs a minimum of X papers to have a Hirsch-Index of
X. Junior and early career researchers will be largely discriminated
against as they may only have a few papers.

\subsubsection{The Hirsch-Index is not helpful to compare researchers
from different
fields}\label{the-hirsch-index-is-not-helpful-to-compare-researchers-from-different-fields}

The strength of the Hirsch-Index, namely its integration of output and
reception and the fact it does not, per definition at least, require
field-normalization, also produces a weakness making it not universally
applicable across different academic disciplines due to variations in
publication and citation practices.

\subsubsection{The Hirsch-Index can be a metric garbage
compactor}\label{the-hirsch-index-can-be-a-metric-garbage-compactor}

As the Hirsch-Index is neither field nor journal normalized it produces
another issue as it ignores almost all context of citations. For
instance, it does not distinguish between citations from high-quality
sources and less reputable ones. In theory, a high Hirsch-Index can be
achieved by a combination of high levels of self-citation and strategic
publication in low impact journals. Or any other means of gaming
citation counts for that matter.

\subsubsection{The Hirsch-Index features strong fluctuations between
databases and
fields}\label{the-hirsch-index-features-strong-fluctuations-between-databases-and-fields}

The Hirsch-Index can vary enormously depending on the database used with
the impact of database selection depending on the field, or, more
precisely the overall publication culture prevalent in a field and the
tendency to publish pre-prints. Usually, differences between Web of
Science (WoS) and Scopus are rather small with large differences being
mostly between the aforementioned and Google Scholar with strong
pre-publishing fields such as computer science receiving very high
premiums in Google Scholar.

\section{Further Reading}\label{further-reading-8}

Bar-Ilan, J. (2008). Which h-index? --- A comparison of WoS, Scopus and
Google Scholar. \emph{Scientometrics, 74}(2), 257--71.
\url{https://doi.org/10.1007/s11192-008-0216-y}

Bornmann, L., \& Daniel, H.-D. (2007). What do we know about the h
index? \emph{Journal of the American Society for Information Science and
Technology, 58}(9), 1381--85. \url{https://doi.org/10.1002/asi.20609}

Egghe, L. (2010). The Hirsch index and related impact measures.
\emph{Annual Review of Information Science and Technology, 44}(1),
65--114. \url{https://doi.org/10.1002/aris.2010.1440440109}

Hirsch, J. E. (2005). An index to quantify an individual's scientific
research output. \emph{Proceedings of the National Academy of Sciences,
102}(46), 16569--72. \url{https://doi.org/10.1073/pnas.0507655102}

Van Raan, A. F. J. (2006). Comparison of the Hirsch-index with standard
bibliometric indicators and with peer judgment for 147 chemistry
research groups. \emph{Scientometrics, 67}(3), 491--502.
\url{https://doi.org/10.1556/Scient.67.2006.3.10}

\chapter{Percentile-Based Indicators (Excellence Rate,
ppTopX)}\label{percentile-based-indicators-excellence-rate-pptopx}

\section{What is it?}\label{what-is-it-1}

In bibliometrics, percentile-based indicators are statistical measures
used to evaluate and compare academic publications. These indicators
place publications within the context of a larger set of data, typically
by showing where a particular work falls in the distribution of all
works in the same field or subject area. The ppTop10 (percentage of
publications in the top 10\%) is a specific type of percentile-based
indicator. It measures the proportion of an entity's (such as an
author's, institution's, or country's) publications that are in the top
10\% of the most cited papers in their respective fields. This is
determined by comparing the citation count of a paper to that of other
papers in the same field and year.

\section{Why is it important?}\label{why-is-it-important-8}

Percentile-based indicators follow two main arguments. The first refers
to the issue of log-normality in the citation distribution, i.e.~a lot
of uncited and rarely cited papers and a very low number of very highly
cited papers leading to a right-skewed distribution, and the discussion
if the arithmetic mean (average) of the number of citations can provide
a useful value in the light of these highly cited outliers. The usual
reply to such problematisations is to use the median as a central value.
In citation distributions the median, by nature being an integer value
rather than a metric number, is often ranging between 0 and 2 (on a good
day), which makes it a very bad metric to capture expectations of
\emph{the usual number of citations in a field}. While the median is
rarely used, the percentile being used is rather the top 10\%, 5\% or
1\%. Which leads to the second argument that also highlights its
colloquial name: The excellence indicator. The argument in this case is
that using the average as a benchmark, e.g.~in field normalization, is
misleading and does not provide substantial and good advice to policy
makers. Rather, focus should be put on the upper parts of the citation
distribution denoting notions of \emph{excellence} rendering everything
else as rather irrelevant \emph{noise}. The share of a unit of
observation in these excellent publications relative to the total amount
produced captures the notion of: ``How much excellence is there relative
to the total output?''

\section{How It Works}\label{how-it-works}

\subsubsection{Step 0: Field and Year
Normalization}\label{step-0-field-and-year-normalization}

Each publication is compared to others within the same field and
publication year. This ensures that the comparison is fair and accounts
for variations in citation practices across different disciplines and
over time. This is crucial as it controls for publication and citation
practices between the different fields.

\subsubsection{Step 1: Ranking the total corpus by citation
count}\label{step-1-ranking-the-total-corpus-by-citation-count}

ALL publications, or more precisely, all citable items, in a field for a
specific publication year in a given database are ranked based on their
citation count.

\subsubsection{Step 2: Identifying Top
10\%}\label{step-2-identifying-top-10}

The top 10\% of publications in terms of citation count are identified.

\subsubsection{Step 3: Calculating
ppTop10}\label{step-3-calculating-pptop10}

The percentage of an entity's publications that fall into this top 10\%
category is calculated. For example, if an institution has 100 papers
and 15 of them are in the top 10\% most cited in their fields, the
institution's ppTop10 score would be 15\%.

A higher ppTop10 score suggests that a significant portion of an
entity's work is highly recognized in its field.

\section{Limitations}\label{limitations-9}

\subsubsection{Excellence rates shift observation and
discussion}\label{excellence-rates-shift-observation-and-discussion}

The strength of ppTop10 also contributes to its major weakness as it
pushed discursive positions toward over-indexing on excellence as the
only meaningful category. This might be useful in some discussions but
shifts the discussion away from discussing emerging fields, topics or
delayed reception. Overall, while ppTop10 indicates high citation
levels, it doesn't account for the overall impact or quality of the
research. Some important works may not be highly cited.

\section{Further Reading}\label{further-reading-9}

Bornmann, L., \& Mutz, R. (2011). Further steps towards an ideal method
of measuring citation performance: The avoidance of citation (ratio)
averages in field-normalization. \emph{Journal of Informetrics, 5}(1),
228--230. \url{https://doi.org/10.1016/j.joi.2010.10.009}

Bornmann, L., \& Williams, R. (2020). An evaluation of percentile
measures of citation impact, and a proposal for making them better.
\emph{Scientometrics, 124}(2), 1457--1478.
\url{https://doi.org/10.1007/s11192-020-03512-7}

Pudovkin, A. I., \& Garfield, E. (2009). Percentile rank and author
superiority indexes for evaluating individual journal articles and the
author's overall citation performance. \emph{Collnet Journal of
Scientometrics and Information Management, 3}(2), 3--10.
\url{https://doi.org/10.1080/09737766.2009.10700871}

Tijssen, R. J. W., Visser, M. S., \& Van Leeuwen, T. N. (2002).
Benchmarking international scientific excellence: Are highly cited
research papers an appropriate frame of reference? \emph{Scientometrics,
54}(3), 381--397. \url{https://doi.org/10.1023/A:1016082432660}

Van Leeuwen, T. N. (2003). The Holy Grail of science policy: Exploring
and combining bibliometric tools in search of scientific excellence.
\emph{Scientometrics, 57}(2), 257--280.
\url{https://doi.org/10.1023/A:1024141819302}

Waltman, L., \& Schreiber, M. (2013). On the calculation of
percentile‐based bibliometric indicators. \emph{Journal of the American
Society for Information Science and Technology, 64}(2), 372--379.
\url{https://doi.org/10.1002/asi.22775}

\chapter{AoR (Averages of Ratios) vs.~RoA (Ratios of Averages) in
Bibliometrics}\label{aor-averages-of-ratios-vs.-roa-ratios-of-averages-in-bibliometrics}

\section{What is it?}\label{what-is-it-2}

Averages (or Means) of Ratios (AoR) and Ratios of Averages (or Means)
(RoA) are two distinct statistical approaches used in bibliometric
analysis and other fields where comparative analysis of data sets is
required. The issue is most prevalent in evaluatory bibliometrics when
using average-based indicators, i.e.~indicators that somehow aim to
normalize output or citations against something to be considered
\emph{the usual}, i.e.~an average. In simple terms, the AoR approach
involves first calculating the ratio for individual pairs (!) of data
points across two data sets, and then finding the average (mean) of
these ratios. This gives equal weight to each data point, focusing on
the individual level of comparison before aggregating. In RoA, by
contrast, this approach calculates the average (mean) for each data set
(!) separately and then takes the ratio of these two means. This method
emphasizes the collective characteristics of each data set, treating the
data sets as aggregated wholes from the onset.

\section{Why is it important?}\label{why-is-it-important-9}

The choice between these two methods can significantly influence the
outcome and interpretation of bibliometric analyses, particularly when
comparing groups or categories. Both methods reflect different aspects
of the data and what to learn from it. AoR can highlight
individual-level differences or similarities, while RoA provides a
comparison at the aggregate level. The selection of the appropriate
method therefore depends on the research question and the nature of the
data, making it crucial to understand both approaches for accurate and
meaningful analysis. AoR typically results in higher scores than RoA,
particularly for entities with fewer papers or lower impact scores. RoA
tends to underestimate the impact, especially for departments,
institutions, and countries with lower scores. RoA eventually leads to
inconsistencies in global averages, which should ideally sum to unity
but don't always do so in practice. The choice between AoR and RoA
impacts the evaluation of research performance, with AoR being generally
more favorable for entities with fewer publications.

\section{How does it work?}\label{how-does-it-work-7}

\subsubsection{AoR:}\label{aor}

\begin{enumerate}
\def\labelenumi{\arabic{enumi}.}
\tightlist
\item
  Calculate the ratio for each corresponding pair of data points across
  the two sets.
\item
  Sum up all these individual ratios.
\item
  Divide the total by the number of ratios to get the mean.
\end{enumerate}

\subsubsection{RoA:}\label{roa}

\begin{enumerate}
\def\labelenumi{\arabic{enumi}.}
\tightlist
\item
  Calculate the mean for each data set separately.
\item
  Divide the mean of one data set by the mean of the other to get the
  ratio.
\end{enumerate}

\section{Limitations}\label{limitations-10}

\subsubsection{AoR:}\label{aor-1}

\begin{itemize}
\tightlist
\item
  Can be heavily influenced by outliers or extreme ratios.
\item
  May not accurately represent the overall characteristics of the data
  sets if there's significant variability within them.
\end{itemize}

\subsubsection{RoA:}\label{roa-1}

\begin{itemize}
\tightlist
\item
  Can mask individual data point variations, leading to oversimplified
  conclusions.
\item
  May not be appropriate in cases where data sets have widely differing
  distributions or sizes.
\end{itemize}

\section{Further Reading}\label{further-reading-10}

Egghe, L. (2012). Averages of ratios compared to ratios of averages:
Mathematical results. \emph{Journal of Informetrics, 6}(2), 307--317.
\url{https://doi.org/10.1016/j.joi.2011.12.007}

Egghe, L., \& Rousseau, R. (1996). Average and global impact of a set of
journals. \emph{Scientometrics, 36}(1), 97--107.
\url{https://doi.org/10.1007/BF02126648}

Larivière, V., \& Gingras, Y. (2011). Averages of ratios vs.~ratios of
averages: An empirical analysis of four levels of aggregation.
\emph{Journal of Informetrics, 5}(3), 392--399.
\url{https://doi.org/10.1016/j.joi.2011.02.001}

Vinkler, P. (1996). Model for quantitative selection of relative
scientometric impact indicators. \emph{Scientometrics, 36}(2), 223--236.
\url{https://doi.org/10.1007/BF02017315}

\part{Explore}

\chapter{Weightedness and Directedness of a Graph in
Bibliometrics}\label{weightedness-and-directedness-of-a-graph-in-bibliometrics}

\section{What is it?}\label{what-is-it-3}

In bibliometrics, and more broadly in network analysis, a graph is a
visual representation of relationships between different entities (like
authors, publications, or institutions). The \emph{weightedness} and
\emph{directedness} of a graph are two important characteristics.
Weightedness refers to whether the edges (connections) in the graph have
weights that signify the strength or intensity of the connection. For
example, in a co-authorship network, the weight could represent the
number of papers co-authored by two researchers. Directedness indicates
whether the edges have a direction, representing asymmetric
relationships. In a citation network, for example, an edge from paper A
to paper B would signify that A cites B, not the other way around.

\section{Why is it Important?}\label{why-is-it-important-10}

Weightedness and directedness are two characteristics that inform what
can be done analysis-wise with network data. For instance, these
characteristics determine what type of community detection algorithm may
be useful and which may produce meaningful results. Directedness refers
to the presence of directed edges in a graph, indicating the
directionality of flows, e.g.~of information, between nodes. In
bibliometrics or altmetrics, directedness can be an important factor
when identifying communities. For example, in a directed social network
where interactions are one-way (e.g., following a user on Twitter),
community detection algorithms may prioritize densely connected groups
of users who receive and share information more frequently with each
other. Weightedness refers to the presence of edges with different
weights in a graph. In the context of community detection, weightedness
can be used to capture the significance or strength of connections
between nodes within a community. For instance, in bibliometrics or
altmetrics, where some users have more followers or higher levels of
influence than others, algorithms may assign higher weights to these
connections and prioritize communities that are densely connected with
influential individuals. Similarly, the analysis of co-authorship
structures on author or institutional level will benefit from
considering the weight (i.e.~the number of co-published papers) rather
than just relying on information if co-authorship occurred in a binary
fashion, i.e.~yes or no. The relevance of directedness and weightedness
in graph analysis becomes evident when considering different algorithms
for community detection. For example, the Louvain method is a popular
community detection algorithm that incorporates both relatedness, via
centrality measures, and weightedness, by using the weights to edges
based on their impact or influence. This approach allows for the
identification of communities with not only high relatedness but also
strong ties to influential nodes within those communities. In a similar
vein, the Leiden algorithm, a rather elegant extension of the Louvain
method, works the same way in this respect. Another community detection
algorithm, the Girvan-Newman algorithm, leverages edge removal to
identify communities based on maximizing betweenness centrality (or
eigenvector centrality for that matter). By considering both
directedness and weightedness, this algorithm can detect densely
connected groups of nodes while also accounting for the significance of
connections within those communities.

Overall, understanding weightedness and directedness is important in
exploratory bibliometric analysis for several reasons. First of all, the
characteristics allow for a more nuanced and accurate representation of
the relationships in a scholarly network and thereby also a more
informed interpretation. Again, it does make a difference if a network
is directed or not, e.g.~if a paper references another paper or is cited
by it. It does make a difference if a network is weighted, i.e.~if a
co-publication between two organizations is framed as simply existing or
not or if a strength (e.g.~the Salton Cosine of co-publications) is
being observed. In short, honoring these characteristics can enable
researchers to conduct more detailed and specific analyses, like
identifying influential authors (in directed graphs) or strong
collaboration networks (in weighted graphs). Furthermore, directed and
weighted graphs provide deeper insights into the nature and intensity of
scholarly interactions, which can be crucial for understanding
scientific communication and collaboration patterns. Weighted networks
sometimes offer clear and informative visual representations of complex
relationships by integrating the strength of ties as width of the
connecting edges. Finally, positionality algorithms will benefit from
weighted graphs as this helps in integrating weights into the
calculator, thereby putting strongly related nodes closer together than
less strongly related nodes.

\section{How Does it Work?}\label{how-does-it-work-8}

Usually the two graph characteristics can be deduced from the way the
network data has been collected and constructed. When using data from a
third party it is always advisable to double check if the assumption
being made, e.g.~an author collaboration matrix being
undirected-weighted, is actually true. In the case of a collaboration
network based on authors researchers may make the decision to give the
relationship between first author (or last author \emph{PI} position
depending on the field) and all the other authors a direction to, for
instance, signify some sort of managerial role. The way such networks
are constructed thereby follows a) the potentiality of what weightedness
and directedness can actually be achieved based on the underlying data,
and b) the research question in mind and the rationale behind the
construction process. Some characteristics can be checked. If weights
greater 1 are present in a dataset, then it is not unweighted. If only
weights of 0 or 1 exist this does NOT imply that the graph is weighted.
It simply can only have relationships of 0 and 1. For all intents and
purposes it CAN be treated as unweighted, though. When cell-wise
subtracting the matrix representing a graph with its transposed form
(rows and columns are switched, i.e.~the matrix is tipped to the side,
and the result for all cells is NOT 0, then the graph should be
directed, assuming no mistake has been made constructing it. If all
cells feature 0, then the graph CAN be undirected. Yet, in the case of
directedness the graph should still be treated as directed.

\section{Limitations}\label{limitations-11}

\subsection{Always check, never assume}\label{always-check-never-assume}

The thing with these characteristics is that sometimes developers of an
algorithm implementation assume that the user knows what they are doing
and take some measures to optimize. For instance, developers might
choose to only use the upper or lower triangle of a matrix representing
a graph as they hold the same information in an undirected graph anyway
and because, assuming N number of nodes, doing \(N*\dfrac{(N-1)}{2}\)
calculations is less than doing \(N*N-N\) calculations. In short: An
implementation can produce results even IF the data does not follow the
requirements of the method.

\chapter{Community Detection Algorithms in
Bibliometrics}\label{community-detection-algorithms-in-bibliometrics}

\section{What is it?}\label{what-is-it-4}

Community detection algorithms in bibliometrics are computational
methods used to identify clusters or groups within scientific networks.
These networks can be based on various relationships such as
co-authorship on author, organization, regional or national level,
citations, co-citations, shared keywords among scientific papers
etc\ldots{} Said clustering can reveal hidden patterns, subfields, or
research trends within a larger scientific domain.

\section{Why is it Important?}\label{why-is-it-important-11}

Visual inspection of networks above a certain size is challenging at
best, especially if networks are unweighted, i.e.~like with the current
way we use citation data, which is largely showing if a paper cites
another paper or not and does not show \emph{how often} a paper mentions
another. Visualizing such networks often leads to the so-called
\emph{hairball problem}. Finding structures in such networks visually(!)
seems an almost impossible task with dense networks. Also relying too
much on attributes of nodes (degree, betweenness\ldots) can usually be
instructive when it comes to identifying \emph{influential nodes}
(authors, organizations), but is of very little use when the question
relates to identifying structures. To achieve this, one way to move
forward is using algorithms that identify \emph{clusters} or
\emph{agglomerations}, in the loosest sense of the word, based mostly on
either using matrix algebra, discrete mathematics, or graph theory.
Common algorithms include modularity optimization methods, hierarchical
clustering, and spectral clustering. Such algorithms come by a plethora
of names: Structure-detection algorithms, community-detection
algorithms, clustering algorithms etc\ldots{} The common denominator is
that these algorithms aim to uncover the structure of these networks by
grouping nodes (again\ldots{} such as authors, papers, or journals) into
communities or clusters that are more densely connected with each other
than with the rest of the network. By doing that, they can facilitate
the analysis of the development, interaction, and interdisciplinary
nature of various scientific fields or help identify emerging research
areas and trends. As a strategic tool, identifying such \emph{clusters}
can guide organization and funding bodies in making informed decisions
about resource allocation, based on the identification of key research
areas and collaboration networks. Finally, using such algorithms can
guide formative evaluation as it can continuously assist researchers in
identifying potential collaborators or relevant research groups.

\section{How Does it Work?}\label{how-does-it-work-9}

Implementing community detection algorithms involves a number of steps
and there are certain flavors how these can be used in detail. Yet,
quite usually, the following steps have to occur.

\subsection{Step 1: Network
Construction}\label{step-1-network-construction}

Building a network from bibliometric data, where nodes represent
entities like authors or papers, and edges represent relationships like
co-authorship or citations.

\subsection{Step 2: Deduce the type of network and select
algorithm}\label{step-2-deduce-the-type-of-network-and-select-algorithm}

Not all algorithms make sense for all kinds of networks. Select an
appropriate algorithm based on the graph being directed or undirected as
well as the graph being weighted or unweighted. Some algorithms require
a fully connected graph not featuring multiple \emph{components},
i.e.~any node can be reached by any other node, or in other words, the
graph has no \emph{un-connected subgraphs}.

\subsection{Step 3: Run analysis}\label{step-3-run-analysis}

Applying community detection algorithms to the network.

\subsection{Step 4 Determine stable number of
clusters/groups}\label{step-4-determine-stable-number-of-clustersgroups}

This usually involves an algorithm that optimizes towards a certain
measure maximizing edge characteristics within clusters, such as average
degree of nodes in a cluster, number of edges within a cluster
etc\ldots{} and minimizing this characteristic between clusters. There
are a good dozen of means to do this (Elbow method, gap statistics, the
Calinski--Harabasz index, the Davies--Bouldin index\ldots)

\subsection{Step 5 Labeling the clusters, Analysis and
Interpretation}\label{step-5-labeling-the-clusters-analysis-and-interpretation}

Analyzing the resulting communities to draw insights about the
underlying structure of the scientific network.

\section{Limitations}\label{limitations-12}

\subsection{Algorithm Complexity}\label{algorithm-complexity}

Some community detection algorithms can be complex and computationally
intensive depending largely on the size and density of the network. The
more sparse and the smaller a network, the faster computation can go.
Also, some algorithms of this sort can be \emph{NP-hard}, i.e.~solvable
in Nondeterministic Polynomial time leading to an exponential increase
in computation relative to increasing the extent (number of nodes and
edges) of the problem at hand making them \emph{tricky} to solve quickly
and easily. (To provide an understanding of what NP-hard means in
practice: Problems being NP-hard is why most non-trivial cryptography
\emph{works}.)

\subsection{Community detection is no silver bullet for
interpretation}\label{community-detection-is-no-silver-bullet-for-interpretation}

The results can sometimes be challenging to interpret, especially in
highly interconnected or multidisciplinary fields or situations of
complex interconnectedness, for instance, when the connectedness is a
phenomenon in and of itself like in the case of interdisciplinarity.

\subsection{Structure Detection algorithms detect
structures}\label{structure-detection-algorithms-detect-structures}

That's what they do. This does not mean that different algorithms
produce the same conclusions or that one conclusion has more
\emph{validity} than another. There are some common metrics that
resemble quality measures from inferential statistics, such as
modularity, but they only are interpretable within very limited
confines. Furthermore, humans are excellent pattern recognition
machines, leading to usually a stance of expressing their own ideas
about structures onto the results, which in turn usually leads more to a
\emph{see what you already know} situation.

\subsection{Dependence on Data
Quality}\label{dependence-on-data-quality}

The effectiveness of community detection is heavily reliant on the
quality and completeness of the bibliometric data.

\section{Further Reading}\label{further-reading-11}

Blondel, V. D., Guillaume, J.-L., Lambiotte, R., \& Lefebvre, E. (2008).
Fast unfolding of communities in large networks. \emph{Journal of
Statistical Mechanics: Theory and Experiment, 2008}(10), P10008.
\url{https://doi.org/10.1088/1742-5468/2008/10/P10008}

Clauset, A., Newman, M. E. J., \& Moore, C. (2004). Finding community
structure in very large networks. \emph{Physical Review E, 70}(6),
066111. \url{https://doi.org/10.1103/PhysRevE.70.066111}

Luke, D. A. (2015). \emph{A user's guide to network analysis in R}. Use
R! Cham: Springer International Publishing.
\url{https://doi.org/10.1007/978-3-319-23883-8}

Newman, M. E. J., \& Girvan, M. (2004). Finding and evaluating community
structure in networks. \emph{Physical Review E, 69}(2), 026113.
\url{https://doi.org/10.1103/PhysRevE.69.026113}

Traag, V. A., Waltman, L., \& Van Eck, N. J. (2019). From Louvain to
Leiden: Guaranteeing well-connected communities. \emph{Scientific
Reports, 9}(1), 5233. \url{https://doi.org/10.1038/s41598-019-41695-z}

Traag, V. A., \& Šubelj, L. (2023). Large network community detection by
fast label propagation. \emph{Scientific Reports, 13}(1), 2701.
\url{https://doi.org/10.1038/s41598-023-29610-z}

\chapter{Bibliographic Coupling}\label{bibliographic-coupling}

\section{What is Bibliographic
Coupling?}\label{what-is-bibliographic-coupling}

Bibliographic coupling (BC) occurs when two scientific papers reference
the same third paper in their bibliographies. In simple terms,
bibliographic coupling is the number of shared references for a pair of
papers. The concept was introduced by M.M. Kessler in 1963 and has
developed into a key tool in bibliometrics for analyzing the structure
and development of scientific research, especially in terms of shared
intellectual heritage contributing to scientific papers suggesting a
direct relationship in subject matter or methodology.

\section{Why is it Important?}\label{why-is-it-important-12}

Bibliographic coupling is important as it helps in identifying research
trends and the development of scientific fields. Papers that share many
references are likely to be closely related, offering insights into how
researchers build upon existing knowledge. This analysis is particularly
useful in identifying research clusters, understanding the evolution of
topics, and detecting emerging areas of study. Also, sometimes it can
help to limit down the size of large networks by using only such
publications that feature a minimum bibliographic coupling of 1.

\section{How Does it Work?}\label{how-does-it-work-10}

Computing bibliographic coupling is rather straightforward. Usually
reference lists from a set of relevant scientific papers are compiled
using a bibliographic database. The preferred method is to produce a
citing paper x cited paper matrix (citation network). This network is an
unweighted-directed graph. Unweighted directed graphs can be transformed
into weighted-undirected graphs (the bibliographic coupling network) by
matrix multiplication using the aforementioned matrix and performing a
matrix multiplication with the transposed matrix. It should be noted
that matrix multiplication, in contrast to its scalar relative, is
sensitive to the order of multiplicand and multiplier. \(A \cdot A^T\)
is not the same as \(A^T \cdot A\). When confusing the order the result
is, interestingly, the co-citation network.

\section{Limitations}\label{limitations-13}

\subsection{Bibliographic coupling is over-indexing on citation
behavior}\label{bibliographic-coupling-is-over-indexing-on-citation-behavior}

Bibliographic coupling is limited to the analysis of reference lists.
Hence, it only reflects the similarities based on the chosen references,
which might miss broader connections such as language use.

\subsection{Bibliographic coupling paints a static
picture}\label{bibliographic-coupling-paints-a-static-picture}

As mentioned above, reference lists are static. In contrast to
co-citation, which can change over time as new publications citing
others are being produced, bibliographic coupling is static, i.e.~it
will not change as the reference lists have already been compiled prior
to final publication. BC therefore offers a snapshot based on the time
of publication, not accounting for subsequent developments.

\subsection{Bibliographic Coupling is biased toward older
contributions}\label{bibliographic-coupling-is-biased-toward-older-contributions}

More recent papers have less opportunity to be bibliographically coupled
due to their shorter presence in the literature.

\subsection{Bibliographic Coupling is dependent on referencing and
citation
culture}\label{bibliographic-coupling-is-dependent-on-referencing-and-citation-culture}

Different academic fields have different referencing practices, which
can affect the outcome of the analysis.

\section{Further Reading}\label{further-reading-12}

Glänzel, W., \& Czerwon, H. J. (1996). A new methodological approach to
bibliographic coupling and its application to the national, regional and
institutional level. \emph{Scientometrics, 37}(2), 195--221.
\url{https://doi.org/10.1007/BF02093621}

Jarneving, B. (2005). A comparison of two bibliometric methods for
mapping of the research front. \emph{Scientometrics, 65}(2), 245--63.
\url{https://doi.org/10.1007/s11192-005-0270-7}

Kessler, M. M. (1963). Bibliographic coupling between scientific papers.
\emph{American Documentation, 14}(1), 10--25.
\url{https://doi.org/10.1002/asi.5090140103}

Kessler, M. M. (1963). Bibliographic coupling extended in time: Ten case
histories. \emph{Information Storage and Retrieval, 1}(4), 169--87.
\url{https://doi.org/10.1016/0020-0271(63)90016-0}

Weinberg, B. H. (1974). Bibliographic coupling: A review.
\emph{Information Storage and Retrieval, 10}(5--6), 189--96.
\url{https://doi.org/10.1016/0020-0271(74)90058-8}

Zhao, D., \& Strotmann, A. (2008). Evolution of research activities and
intellectual influences in information science 1996--2005: Introducing
author bibliographic-coupling analysis. \emph{Journal of the American
Society for Information Science and Technology, 59}(13), 2070--86.
\url{https://doi.org/10.1002/asi.20910}

\chapter{Co-Citation Analysis}\label{co-citation-analysis}

\section{What is co-citation
analysis?}\label{what-is-co-citation-analysis}

Co-citation analysis is a method used in bibliometrics to study the
relationship between scientific publications and, in some of the
derivatives, such as author-co-citation analysis (ACA), between other
entities (nodes) that can be attributed to scientific publications. Two
documents are co-cited, when they both receive (!) a citation from the
same later (!) document. The strength of co-citation equals the number
of citing documents for which this is the case, i.e.~a co-citation of
strength 3 means that 3 documents cite both scientific publications
under analysis. This implies that an analysis of co-citation can be
performed for each combination of citable items in a bibliographic
dataset. Essentially, if two papers are often cited together, it implies
a topical or methodological relationship between them that in turn
implies intellectual heritage. This concept was introduced by Henry
Small in the 1970s and is based on his notion that citations are
\emph{concept symbols} bringing together ideas from bibliometrics and
semiotics, the latter being quite fashionable at that time. It shares a
number of characteristics with bibliographic coupling (shared entries in
reference lists between different papers), a concept that also relates
to notions of \emph{intellectual heritage}.

\section{Why is co-citation analysis
important?}\label{why-is-co-citation-analysis-important}

Co-citation analysis is important because it helps to map the structure,
evolution and dynamics of scientific fields as they can change over
time, with new scientific publications being produced that can
contribute to co-citation structures. By \emph{being attributed
together}, co-citation reveals how knowledge areas are interconnected
and can identify emerging trends and, with sufficient data and
computational power, even the emergence of (sub-)disciplines. The
analysis is also crucial for understanding the influence of a particular
work in a field, as frequent co-citations indicate a shared relevance or
foundational and conceptual status among the cited works.

\section{How Does it Work?}\label{how-does-it-work-11}

See Bibliographic Coupling but reverse the order of multiplicand and
multiplier when performing matrix multiplication.

\section{Limitations}\label{limitations-14}

\subsection{Calculating Co-Citation networks can be computationally
intensive}\label{calculating-co-citation-networks-can-be-computationally-intensive}

Citation networks can grow large very quickly, which is mainly due to
the dominant nature of most phenomena in bibliometrics being
\emph{log-normal} distributed. The probably most efficient one-shot
computational method to calculate co-citation is matrix multiplication.
Yet, computing matrix multiplication is currently considered
\emph{P-hard}, which means that even though it actually can be solved in
polynomial time, the computational time required for matrix
multiplication grows rapidly, making it an expensive operation for large
matrices, which in turn means that even if the number of papers to
observe might be rather small, the dispersedness of citations can
produce large matrices\ldots{} and hence produce long computational
times.

\subsection{Calculating co-citation can require substantial
computational
resources}\label{calculating-co-citation-can-require-substantial-computational-resources}

As stated above, co-citation can be data intensive, which means that,
despite the challenge of computation it also faces the challenge of very
high quality data. In short: While the point about computational time is
a \emph{CPU (or GPU) problem}, the issue here is a \emph{RAM (memory)
problem}. The problems are not the same. Parallelizing the problem to
decrease computation time, might come to a grinding halt because RAM is
obliterated by the parallel computing threads and the large amounts of
data.

\subsection{Co-Citation is less immediate compared to bibliographic
coupling}\label{co-citation-is-less-immediate-compared-to-bibliographic-coupling}

In bibliographic coupling, older publications might be over-represented.
In co-citation the opposite is true as there is a time-lag involved (see
citation windows). Citation data may not reflect the most current trends
due to publication and citation delays. Yet, citation windows usually
are not being applied (in contrast to evaluative bibliometrics).

\subsection{Co-Citation is dependent on citation
practices}\label{co-citation-is-dependent-on-citation-practices}

Variations in citation practices across fields can affect the results.
It may therefore over- or under-represent some structures based in
inter- or transdisciplinary fields.

\subsection{\texorpdfstring{Co-citation has a bias towards
\emph{disciplinarity}}{Co-citation has a bias towards disciplinarity}}\label{co-citation-has-a-bias-towards-disciplinarity}

Even though co-citation can be used to identify new topics and fields,
there is a bias involved that relates to the point about citation
practices. Highly interdisciplinary work might be underrepresented in
co-citation networks if it doesn't fit neatly into established citation
patterns.

\section{Further Reading}\label{further-reading-13}

Small, H. (1973). Co-citation in the scientific literature: A new
measure of the relationship between two documents. \emph{Journal of the
American Society for Information Science, 24}(4), 265--69.
\url{https://doi.org/10.1002/asi.4630240406}

White, H. D., \& Griffith, B. C. (1981). Author cocitation: A literature
measure of intellectual structure. \emph{Journal of the American Society
for Information Science, 32}(3), 163--71.
\url{https://doi.org/10.1002/asi.4630320302}

\bookmarksetup{startatroot}

\chapter{Bibliography}\label{bibliography}

Aksnes, D. W., Schneider, J. W., \& Gunnarsson, M. (2012). Ranking
national research systems by citation indicators. A comparative analysis
using whole and fractionalised counting methods. \emph{Journal of
Informetrics, 6}(1), 36--43.
\url{https://doi.org/10.1016/j.joi.2011.08.002}

Baeza-Yates, R., \& Ribeiro-Neto, B. (2011). \emph{Modern information
retrieval: The concepts and technology behind search} (2nd ed.).
Addison-Wesley Publishing Company.

Bailón-Moreno, R., Jurado-Alameda, E., Ruiz-Baños, R., \& Courtial, J.
P. (2005). Bibliometric laws: Empirical flaws of fit.
\emph{Scientometrics, 63}(2), 209--229.
\url{https://doi.org/10.1007/s11192-005-0211-5}

Bar-Ilan, J. (2008). Which h-index? --- A comparison of WoS, Scopus and
Google Scholar. \emph{Scientometrics, 74}(2), 257--271.
\url{https://doi.org/10.1007/s11192-008-0216-y}

Blondel, V. D., Guillaume, J.-L., Lambiotte, R., \& Lefebvre, E. (2008).
Fast unfolding of communities in large networks. \emph{Journal of
Statistical Mechanics: Theory and Experiment, 2008}(10), P10008.
\url{https://doi.org/10.1088/1742-5468/2008/10/P10008}

Bornmann, L. (2007). How can citation impact in bibliometrics be
normalized? A new approach combining citing-side normalization and
citation percentiles. \emph{Quantitative Science Studies, 1}(4),
1553--1569. \url{https://doi.org/10.1162/qss_a_00089}

Bornmann, L. (2020). An evaluation of percentile measures of citation
impact, and a proposal for making them better. \emph{Scientometrics,
124}(2), 1457--1478. \url{https://doi.org/10.1007/s11192-020-03512-7}

Bornmann, L., \& Daniel, H.-D. (2007). What do we know about the h
index? \emph{Journal of the American Society for Information Science and
Technology, 58}(9), 1381--1385. \url{https://doi.org/10.1002/asi.20609}

Bornmann, L., \& Mutz, R. (2011). Further steps towards an ideal method
of measuring citation performance: The avoidance of citation (ratio)
averages in field-normalization. \emph{Journal of Informetrics, 5}(1),
228--230. \url{https://doi.org/10.1016/j.joi.2010.10.009}

Bornmann, L., \& Williams, R. (2020). An evaluation of percentile
measures of citation impact, and a proposal for making them better.
\emph{Scientometrics, 124}(2), 1457--1478.
\url{https://doi.org/10.1007/s11192-020-03512-7}

Bornmann, L., Haunschild, R., \& Mutz, R. (2020). Should citations be
field-normalized in evaluative bibliometrics? An empirical analysis
based on propensity score matching. \emph{Journal of Informetrics,
14}(4), 101098. \url{https://doi.org/10.1016/j.joi.2020.101098}

Bradford, S. C. (1934). Sources of information on specific subjects.
\emph{Engineering, 26}(4), 85--86.

Campanario, J. M. (2011). Empirical study of journal impact factors
obtained using the classical two-year citation window versus a five-year
citation window. \emph{Scientometrics, 87}(1), 189--204.
\url{https://doi.org/10.1007/s11192-010-0334-1}

Clauset, A., Newman, M. E. J., \& Moore, C. (2004). Finding community
structure in very large networks. \emph{Physical Review E, 70}(6),
066111. \url{https://doi.org/10.1103/PhysRevE.70.066111}

Cleverdon, C. W. (1972). On the inverse relationship of recall and
precision. \emph{Journal of Documentation, 28}(3), 195--201.
\url{https://doi.org/10.1108/eb026538}

Cronin, B. (2001). Hyperauthorship: A postmodern perversion or evidence
of a structural shift in scholarly communication practices?
\emph{Journal of the American Society for Information Science and
Technology, 52}(7), 558--569. \url{https://doi.org/10.1002/asi.1097}

Daniel, B. Klein, \& Chiang, E. (2004). The social science citation
index: A black box---with an ideological bias? \emph{Econ Journal Watch,
1}(1), 134--165.

Daraio, C., Lenzerini, M., Leporelli, C., Naggar, P., Bonaccorsi, A., \&
Bartolucci, A. (2016). The advantages of an ontology-based data
management approach: Openness, interoperability and data quality.
\emph{Scientometrics, 108}(1), 441--455.
\url{https://doi.org/10.1007/s11192-016-1913-6}

Donner, P. (2018). Effect of publication month on citation impact.
\emph{Journal of Informetrics, 12}(1), 330--343.
\url{https://doi.org/10.1016/j.joi.2018.01.012}

Donner, P. (2024). Remarks on modified fractional counting.
\emph{Journal of Informetrics, 18}(4), 101585.
\url{https://doi.org/10.1016/j.joi.2024.101585}

Donner, P., Rimmert, C., \& Van Eck, N. J. (2020). Comparing
institutional-level bibliometric research performance indicator values
based on different affiliation disambiguation systems.
\emph{Quantitative Science Studies, 1}(1), 150--170.
\url{https://doi.org/10.1162/qss_a_00013}

Egghe, L. (2010). The Hirsch index and related impact measures.
\emph{Annual Review of Information Science and Technology, 44}(1),
65--114. \url{https://doi.org/10.1002/aris.2010.1440440109}

Egghe, L. (2012). Averages of ratios compared to ratios of averages:
Mathematical results. \emph{Journal of Informetrics, 6}(2), 307--317.
\url{https://doi.org/10.1016/j.joi.2011.12.007}

Egghe, L., \& Rousseau, R. (1996). Average and global impact of a set of
journals. \emph{Scientometrics, 36}(1), 97--107.
\url{https://doi.org/10.1007/BF02126648}

El Gibari, S., Gómez, T., \& Ruiz, F. (2022). Combining reference point
based composite indicators with data envelopment analysis: Application
to the assessment of universities. \emph{Scientometrics, 127}(8),
4363--4395. \url{https://doi.org/10.1007/s11192-022-04436-0}

Gauffriau, M., Larsen, P. O., Maye, I., Roulin-Perriard, A., \& Von Ins,
M. (2008). Comparisons of results of publication counting using
different methods. \emph{Scientometrics, 77}(1), 147--176.
\url{https://doi.org/10.1007/s11192-007-1934-2}

Glänzel, W. (2004). Towards a model for diachronous and synchronous
citation analyses. \emph{Scientometrics, 60}(3), 511--522.
\url{https://doi.org/10.1023/B:SCIE.0000034391.06240.2a}

Glänzel, W., \& Czerwon, H. J. (1996). A new methodological approach to
bibliographic coupling and its application to the national, regional and
institutional level. \emph{Scientometrics, 37}(2), 195--221.
\url{https://doi.org/10.1007/BF02093621}

Glänzel, W., Schlemmer, B., \& Thijs, B. (2003). Better late than never?
On the chance to become highly cited only beyond the standard
bibliometric time horizon. \emph{Scientometrics, 58}(3), 571--586.
\url{https://doi.org/10.1023/B:SCIE.0000006881.30700.ea}

Glänzel, W., \& Thijs, B. (2012). Using `core documents' for detecting
and labelling new emerging topics. \emph{Scientometrics, 91}(2),
399--416. \url{https://doi.org/10.1007/s11192-011-0591-7}

Harzing, A.-W., \& Alakangas, S. (2016). Google Scholar, Scopus and the
Web of Science: A longitudinal and cross-disciplinary comparison.
\emph{Scientometrics, 106}(2), 787--804.
\url{https://doi.org/10.1007/s11192-015-1798-9}

Heintz, B. (2010). Numerische Differenz. Überlegungen zu einer
Soziologie des (quantitativen) Vergleichs / Numerical difference. Toward
a sociology of (quantitative) comparisons. \emph{Zeitschrift für
Soziologie, 39}(3), 162--181.
\url{https://doi.org/10.1515/zfsoz-2010-0301}

Hicks, D., Wouters, P., Waltman, L., de Rijcke, S., \& Rafols, I.
(2015). Bibliometrics: The Leiden Manifesto for Research Metrics.
\emph{Nature, 520}(7548), 429--431.
\url{https://doi.org/10.1038/520429a}

Hirsch, J. E. (2005). An index to quantify an individual's scientific
research output. \emph{Proceedings of the National Academy of Sciences,
102}(46), 16569--16572. \url{https://doi.org/10.1073/pnas.0507655102}

Jarneving, B. (2005). A comparison of two bibliometric methods for
mapping of the research front. \emph{Scientometrics, 65}(2), 245--263.
\url{https://doi.org/10.1007/s11192-005-0270-7}

Johnes, J. (2018). University rankings: What do they really show?
\emph{Scientometrics, 115}(1), 585--606.
\url{https://doi.org/10.1007/s11192-018-2666-1}

Kessler, M. M. (1963). Bibliographic coupling between scientific papers.
\emph{American Documentation, 14}(1), 10--25.
\url{https://doi.org/10.1002/asi.5090140103}

Kessler, M. M. (1963). Bibliographic coupling extended in time: Ten case
histories. \emph{Information Storage and Retrieval, 1}(4), 169--187.
\url{https://doi.org/10.1016/0020-0271(63)90016-0}

Larivière, V., \& Gingras, Y. (2011). Averages of ratios vs.~ratios of
averages: An empirical analysis of four levels of aggregation.
\emph{Journal of Informetrics, 5}(3), 392--399.
\url{https://doi.org/10.1016/j.joi.2011.02.001}

Larsen, P. O., \& Von Ins, M. (2010). The rate of growth in scientific
publication and the decline in coverage provided by Science Citation
Index. \emph{Scientometrics, 84}(3), 575--603.
\url{https://doi.org/10.1007/s11192-010-0202-z}

Leydesdorff, L., \& Opthof, T. (2010). Normalization at the field level:
Fractional counting of citations.
\url{https://doi.org/10.48550/ARXIV.1006.2896}

Luke, D. A. (2015). \emph{A user's guide to network analysis in R}. Use
R! Cham: Springer International Publishing.
\url{https://doi.org/10.1007/978-3-319-23883-8}

Makkonen, T., \& Van Der Have, R. P. (2013). Benchmarking regional
innovative performance: Composite measures and direct innovation counts.
\emph{Scientometrics, 94}(1), 247--262.
\url{https://doi.org/10.1007/s11192-012-0753-2}

Manning, C. D., Raghavan, P., \& Schütze, H. (2009). \emph{Introduction
to information retrieval}. Cambridge University Press.
\url{https://nlp.stanford.edu/IR-book/pdf/irbookonlinereading.pdf}

Merton, R. K. (1968). The Matthew effect in science: The reward and
communication systems of science are considered. \emph{Science,
159}(3810), 56--63. \url{https://doi.org/10.1126/science.159.3810.56}

Moed, H. F. (2010). Measuring contextual citation impact of scientific
journals. \emph{Journal of Informetrics, 4}(3), 265--277.
\url{https://doi.org/10.1016/j.joi.2010.01.002}

Mongeon, P., \& Paul-Hus, A. (2016). The journal coverage of Web of
Science and Scopus: A comparative analysis. \emph{Scientometrics,
106}(1), 213--228. \url{https://doi.org/10.1007/s11192-015-1765-5}

Moon, H. S., \& Lee, J. D. (2005). A fuzzy set theory approach to
national composite S\&T indices. \emph{Scientometrics, 64}(1), 67--83.
\url{https://doi.org/10.1007/s11192-005-0238-7}

Müller, M. C., Reitz, F., \& Roy, N. (2017). Data sets for author name
disambiguation: An empirical analysis and a new resource.
\emph{Scientometrics, 111}(3), 1467--1500.
\url{https://doi.org/10.1007/s11192-017-2363-5}

Mutschke, P., \& Mayr, P. (2015). Science models for search: A study on
combining scholarly information retrieval and scientometrics.
\emph{Scientometrics, 102}(3), 2323--2345.
\url{https://doi.org/10.1007/s11192-014-1485-2}

Nasir, A., Ali, T. M., Shahdin, S., \& Rahman, T. U. (2011). Technology
achievement index 2009: Ranking and comparative study of nations.
\emph{Scientometrics, 87}(1), 41--62.
\url{https://doi.org/10.1007/s11192-010-0285-6}

Nicolaisen, J., \& Hjørland, B. (2007). Practical potentials of
Bradford's law: A critical examination of the received view.
\emph{Journal of Documentation, 63}(3), 359--377.
\url{https://doi.org/10.1108/00220410710743298}

OECD/European Union/EC-JRC. (2008). \emph{Handbook on constructing
composite indicators: Methodology and user guide}. Paris: OECD
Publishing. \url{https://doi.org/10.1787/9789264043466-en}

Powers, D. M. W. (2008). Evaluation: From precision, recall and
F-Measure to ROC, informedness, markedness and correlation. \emph{Mach.
Learn. Technol., 2}. \url{https://doi.org/10.48550/ARXIV.2010.16061}

Pudovkin, A. I., \& Garfield, E. (2009). Percentile rank and author
superiority indexes for evaluating individual journal articles and the
author's overall citation performance. \emph{Collnet Journal of
Scientometrics and Information Management, 3}(2), 3--10.
\url{https://doi.org/10.1080/09737766.2009.10700871}

Rimmert, C., Schwechheimer, H., \& Winterhager, M. (2017).
Disambiguation of author addresses in bibliometric databases - technical
report. Bielefeld: Universität Bielefeld, Institute for
Interdisciplinary Studies of Science (I²SoS).

Shenton, A. K., \& Hay-Gibson, N. V. (2011). Bradford's Law and its
relevance to researchers. \emph{Education for Information, 27}(4),
217--230. \url{https://doi.org/10.3233/EFI-2009-0882}

Sivertsen, G., Rousseau, R., \& Zhang, L. (2019). Measuring scientific
contributions with modified fractional counting. \emph{Journal of
Informetrics, 13}(2), 679--694.
\url{https://doi.org/10.1016/j.joi.2019.03.010}

Small, H. (1973). Co‐citation in the scientific literature: A new
measure of the relationship between two documents. \emph{Journal of the
American Society for Information Science, 24}(4), 265--269.
\url{https://doi.org/10.1002/asi.4630240406}

Stahlschmidt, S., \& Stephen, D. (2020). Comparison of Web of Science,
Scopus and Dimensions databases. Berlin: Deutsches Zentrum für
Hochschul- und Wissenschaftsforschung.
\url{https://bibliometrie.info/downloads/DZHW-Comparison-DIM-SCP-WOS.PDF}

Tang, L., \& Walsh, J. P. (2010). Bibliometric fingerprints: Name
disambiguation based on approximate structure equivalence of cognitive
maps. \emph{Scientometrics, 84}(3), 763--784.
\url{https://doi.org/10.1007/s11192-010-0196-6}

Tijssen, R. J. W., Visser, M. S., \& Van Leeuwen, T. N. (2002).
Benchmarking international scientific excellence: Are highly cited
research papers an appropriate frame of reference? \emph{Scientometrics,
54}(3), 381--397. \url{https://doi.org/10.1023/A:1016082432660}

Traag, V. A., Waltman, L., \& Van Eck, N. J. (2019). From Louvain to
Leiden: Guaranteeing well-connected communities. \emph{Scientific
Reports, 9}(1), 5233. \url{https://doi.org/10.1038/s41598-019-41695-z}

Traag, V. A., \& Šubelj, L. (2023). Large network community detection by
fast label propagation. \emph{Scientific Reports, 13}(1), 2701.
\url{https://doi.org/10.1038/s41598-023-29610-z}

Van Raan, A. F. J. (2006). Comparison of the Hirsch-index with standard
bibliometric indicators and with peer judgment for 147 chemistry
research groups. \emph{Scientometrics, 67}(3), 491--502.
\url{https://doi.org/10.1556/Scient.67.2006.3.10}

Vinkler, P. (2006). Composite scientometric indicators for evaluating
publications of research institutes. \emph{Scientometrics, 68}(3),
629--642. \url{https://doi.org/10.1007/s11192-006-0123-z}

Waltman, L., \& Van Eck, N. J. (2015). Field-normalized citation impact
indicators and the choice of an appropriate counting method.
\emph{Journal of Informetrics, 9}(4), 872--894.
\url{https://doi.org/10.1016/j.joi.2015.08.001}

Wang, J. (2013). Citation time window choice for research impact
evaluation. \emph{Scientometrics, 94}(3), 851--872.
\url{https://doi.org/10.1007/s11192-012-0775-9}

Weinberg, B. H. (1974). Bibliographic coupling: A review.
\emph{Information Storage and Retrieval, 10}(5--6), 189--196.
\url{https://doi.org/10.1016/0020-0271(74)90058-8}

White, H. D., \& Griffith, B. C. (1981). Author cocitation: A literature
measure of intellectual structure. \emph{Journal of the American Society
for Information Science, 32}(3), 163--171.
\url{https://doi.org/10.1002/asi.4630320302}

Zhao, D., \& Strotmann, A. (2008). Evolution of research activities and
intellectual influences in information science 1996--2005: Introducing
author bibliographic‐coupling analysis. \emph{Journal of the American
Society for Information Science and Technology, 59}(13), 2070--2086.
\url{https://doi.org/10.1002/asi.20910}




\end{document}
